% !TeX program = xelatex
\documentclass[12pt, a4paper]{article}

\usepackage[margin=1in]{geometry}
\usepackage[utf8]{inputenc}
\usepackage[T1]{fontenc}
\usepackage{amsmath, amssymb, amsfonts}
\usepackage{array}
\usepackage{booktabs}
\usepackage{longtable}
\usepackage{parskip}
% The simulation-algorithm box spans more than one page, so it uses the
% page-breakable framed environment rather than a fixed-height minipage.
\usepackage{framed}
\usepackage{setspace}
\usepackage{microtype}
\usepackage[hidelinks]{hyperref}
\onehalfspacing

% Overleaf can retain output.aux from an earlier BibLaTeX build. These no-op
% definitions allow one successful run to ignore obsolete BibLaTeX bookkeeping
% commands and overwrite the auxiliary file for this non-bibliographic document.
\makeatletter
\providecommand{\abx@aux@refcontext}[1]{}
\providecommand{\abx@aux@cite}[2]{}
\providecommand{\abx@aux@segm}[3]{}
\providecommand{\abx@aux@page}[2]{}
\providecommand{\abx@aux@fnpage}[2]{}
\providecommand{\abx@aux@count}[2]{}
\providecommand{\abx@aux@fncount}[2]{}
\providecommand{\abx@aux@backref}[5]{}
\providecommand{\abx@aux@refsection}[2]{}
\providecommand{\abx@aux@biblist}[1]{}
\providecommand{\abx@aux@nociteall}{}
\providecommand{\abx@aux@category}[2]{}
\providecommand{\abx@aux@defaultrefcontext}[3]{}
\providecommand{\abx@aux@defaultlabelprefix}[3]{}
\providecommand{\abx@aux@read@bblrerun}{}
\providecommand{\abx@aux@read@bbl@mdfivesum}[1]{}
\providecommand{\abx@aux@number}[5]{}
\makeatother

% --- Define operators for consistent spacing/typography ---
\DeclareMathOperator{\clip}{clip}
\DeclareMathOperator{\Corr}{Corr}
\DeclareMathOperator{\Std}{Std}
\DeclareMathOperator{\MA}{MA}
\DeclareMathOperator{\RSI}{RSI}
\newcommand{\pospart}[1]{\left(#1\right)^{+}}
\newcommand{\indicator}{\mathbf{1}}
\newcolumntype{L}[1]{>{\raggedright\arraybackslash}p{#1}}
\title{Formalism of the Behavioural Agent-Based Market Model}
\author{}
\date{}
\begin{document}
    \maketitle

    \begin{abstract}
        This document formalises the current C++ implementation of the behavioural
        agent-based market model used in the PhD project.
    \end{abstract}

    \tableofcontents
    \clearpage

    \section{Purpose and Implementation Scope}\label{sec:purpose_scope}

    The model is designed to study whether heterogeneous behavioural agents can
    generate empirical market regularities such as fat-tailed returns, weak
    linear-return autocorrelation, and volatility clustering during crisis-like
    market conditions. The formalism therefore records both the economic mechanism
    and the implementation convention required for reproducible experiments.

    The analysed trading phase uses a persistent continuous double-auction limit
    order book. Fundamental-value shocks, technical trading rules, Prospect Theory
    preferences, social opinions, and noise-trader beliefs affect prices only
    through submitted orders and executions. The official market price changes when
    an endogenous transaction occurs; otherwise it is carried forward. Warm-up
    observations are used only to initialise price history for technical indicators
    and realised-volatility estimates before the analysed simulation phase begins.

    The market contains \(n\ge 1\) tradable risky assets. Every asset follows the
    same theory: a geometric fundamental-value process, one persistent order book
    of its own, the same behavioural rules, and the same stylised-fact screening.
    Assets differ only in their \emph{parameters} --- volatility, drift, price
    level, tick size, and share supply --- never in functional form, so no asset
    is analytically privileged and cross-asset comparisons of returns are
    like-for-like. Fundamental innovations are correlated across assets, which is
    what limits the diversification genuinely available to agents, and agents
    allocate across all assets jointly through a full Markowitz rule using a
    rolling log-return covariance matrix. Setting \(n=1\), or disabling all but
    one declared asset, collapses the model exactly onto the single-asset
    specification. One asset may be nominated \emph{primary}; this nomination
    governs only the legacy single-asset projection of the outputs and carries no
    analytical privilege.

    \section*{Nomenclature}
    \addcontentsline{toc}{section}{Nomenclature}

    \subsubsection*{Sets and Indices}

    \begin{longtable}{@{}L{0.22\textwidth}L{0.50\textwidth}L{0.20\textwidth}@{}}
        \toprule
        \textbf{Symbol} & \textbf{Meaning} & \textbf{Unit/Domain} \\
        \midrule
        \endfirsthead
        \toprule
        \textbf{Symbol} & \textbf{Meaning} & \textbf{Unit/Domain} \\
        \midrule
        \endhead
        \(\mathcal{A}\) & Set of all agents. & Finite set \\
        \(N=|\mathcal{A}|\) & Total number of agents. & Positive integer \\
        \(\mathcal{A}^F\subset\mathcal{A}\) & Set of fundamentalist agents. & Subset of agents \\
        \(\mathcal{A}^C\subset\mathcal{A}\) & Set of chartist agents. & Subset of agents \\
        \(\mathcal{A}^N\subset\mathcal{A}\) & Set of noise-trader agents. & Subset of agents \\
        \(\mathcal{A}^H\subset\mathcal{A}\) & Set of passive buy-and-hold agents. & Subset of agents \\
        \(i,l\in\mathcal{A}\) & Agent indices. & Agent identifiers \\
        \(t\) & Discrete simulation time index. & Non-negative integer \\
        \(\mathcal{J}\) & Set of enabled tradable risky assets, in declared index order. & Finite ordered set \\
        \(n=|\mathcal{J}|\) & Number of enabled assets. & Positive integer \\
        \(j,k\in\mathcal{J}\) & Asset indices. & Asset identifiers \\
        \(j^{\ast}\) & Primary asset index; selects the legacy single-asset projection of the outputs only. & Asset identifier \\
        \(\mathcal{J}^{0}\) & Assets declared but disabled; retained for provenance and given no book, weight, or output row. & Finite set \\
        \bottomrule
    \end{longtable}

    Agent indices use \(i\) and \(l\); the symbol \(j\) is reserved throughout for
    assets. Where a quantity exists once per asset it carries the asset subscript
    \(j\); the single-asset model is recovered by fixing \(n=1\) and dropping that
    subscript.

    \subsubsection*{Prices and State Variables}

    \begin{longtable}{@{}L{0.22\textwidth}L{0.50\textwidth}L{0.20\textwidth}@{}}
        \toprule
        \textbf{Symbol} & \textbf{Meaning} & \textbf{Unit/Domain} \\
        \midrule
        \endfirsthead
        \toprule
        \textbf{Symbol} & \textbf{Meaning} & \textbf{Unit/Domain} \\
        \midrule
        \endhead
        \(F_{j,t}\) & Latent fundamental value of asset \(j\) at time \(t\). & Price units \\
        \(P_{j,t}\) & Observed market price of asset \(j\) at time \(t\). & Price units, \(P_{j,t}>0\) \\
        \(m_{j,t}=\log(F_{j,t}/P_{j,t})\) & Log mispricing of asset \(j\); positive values indicate undervaluation. & Log-price gap \\
        \(c_{i,t}\) & Cash holdings of agent \(i\) at time \(t\); one shared cash account across assets. & Currency units \\
        \(c_0\) & Configured initial cash endowment per agent. & Currency units \\
        \(x_{i,j,t}\) & Holdings of asset \(j\) by agent \(i\) at time \(t\). & Shares or units \\
        \(x_{j,0}\) & Configured initial endowment of asset \(j\) per agent. & Shares or units \\
        \(W_{i,t}=c_{i,t}+\sum_{j}x_{i,j,t}P_{j,t}\) & Mark-to-market wealth of agent \(i\) across all assets. & Currency units \\
        \(\Gamma_{i,t}\) & Cumulative transaction costs paid by agent \(i\). & Currency units \\
        \(\Theta_t\) & Total transaction costs paid by all agents in period \(t\). & Currency units \\
        \(s_{i,j,t}^{-}=\pospart{-x_{i,j,t}}\) & Absolute short position of agent \(i\) in asset \(j\). & Shares or units \\
        \(\mathrm{SI}_{j,t}=\sum_i s_{i,j,t}^{-}\) & Aggregate short interest in asset \(j\) before new orders. & Shares or units \\
        \(\bar X_{j,0}\) & Initial share supply of asset \(j\), used as its securities-lending base. & Shares or units \\
        \(\bar w_{j,t}\) & Cross-sectional mean portfolio weight held in asset \(j\) after the step. & Dimensionless weight \\
        \bottomrule
    \end{longtable}

    \subsubsection*{Trading and Strategy Variables}

    \begin{longtable}{@{}L{0.22\textwidth}L{0.50\textwidth}L{0.20\textwidth}@{}}
        \toprule
        \textbf{Symbol} & \textbf{Meaning} & \textbf{Unit/Domain} \\
        \midrule
        \endfirsthead
        \toprule
        \textbf{Symbol} & \textbf{Meaning} & \textbf{Unit/Domain} \\
        \midrule
        \endhead
        \(a_{i,t}\) & Activity indicator for agent \(i\) at time \(t\). & \(\{0,1\}\) \\
        \(\alpha\) & Probability that an agent is active in a period. & \([0,1]\) \\
        \(\mu_{i,j,t}\) & Expected return on asset \(j\) used in portfolio choice. & Per-period log return \\
        \(\hat r_{i,j,t}^{C},\hat r_{i,j,t}^{F}\) & Baseline expected excess return of chartist and fundamentalist agents on asset \(j\). & Per-period excess log return \\
        \(\breve r_{i,j,t}^{C},\breve r_{i,j,t}^{F}\) & Prospect-adjusted excess return used in allocation when CPT is enabled. & Per-period excess log return \\
        \(\breve{\mathbf r}_{i,t}\) & Vector of prospect-adjusted excess returns across assets. & \(\mathbb{R}^{n}\) \\
        \(\Sigma_t\) & Rolling sample covariance matrix of one-period log returns, shared by all agents. & \(\mathbb{R}^{n\times n}\), symmetric \\
        \(\sigma_{j,t}=\sqrt{\Sigma_{t,jj}}\) & Rolling volatility of asset \(j\) taken from the covariance diagonal. & Return standard deviation \\
        \(w_{i,j,t}\) & Target portfolio weight of agent \(i\) in asset \(j\). & Dimensionless weight \\
        \(\underline w_{i,j},\overline w_{i,j}\) & Effective lower and upper portfolio-weight bounds applied to every asset. & Dimensionless bounds \\
        \(\underline w_i^{\text{cfg}},\overline w_i^{\text{cfg}}\) & Configured allocation floor and ceiling of agent \(i\)'s strategy class. & Dimensionless bounds \\
        \(\bar w_i^{\text{short}}\) & Chartist short-exposure cap; binds jointly with the configured floor. & \([0,1]\) \\
        \(G_i\) & Cap on gross exposure \(\sum_j|w_{i,j,t}|\); non-positive values disable the cap. & Positive real number \\
        \(x_{i,j,t}^{\ast}\) & Target position in asset \(j\). & Shares or units \\
        \(x_{i,j,t}^{\ast,F}\) & Fundamentalist target position in asset \(j\). & Shares or units \\
        \(\mathrm{side}_{i,j,t}\) & Submitted order side of agent \(i\) in asset \(j\). & \(\{\mathrm{Buy},\mathrm{Sell}\}\) \\
        \(q_{i,j,t}\) & Submitted order quantity of agent \(i\) in asset \(j\). & Shares or units, \(q_{i,j,t}\ge0\) \\
        \(\bar q_{i,j,t}\) & Feasible order quantity after cash, inventory, margin, and borrow-availability checks. & Shares or units \\
        \(q_{i,j,t}^{C},q_{i,j,t}^{F},q_{i,j,t}^{N},q_{i,j,t}^{H}\) & Submitted order quantity of chartist, fundamentalist, noise-trader, and buy-and-hold agents. & Shares or units \\
        \(\tilde c_{i,t}\) & Running order-generation cash budget, decremented as buy orders are queued across assets. & Currency units \\
        \(R_{i,t}^{B}\) & Reserved new short-sale shares exposed by the risk-control helper; the current LOB execution path passes zero because borrow capacity is revalidated and short interest is updated at execution time. & Shares or units \\
        \(b_{i,j,t}^{\text{cover}}\) & Forced buy-in quantity used to reduce an under-margined short position in asset \(j\). & Shares or units \\
        \(AS_{j,t}\) & Realised gross allocative surplus from asset \(j\)'s endogenous continuous-double-auction executions in period \(t\). & Currency units \\
        \(MS_{j,t}\) & Maximum feasible gross allocative surplus available from asset \(j\)'s surviving book liquidity and new submitted orders. & Currency units \\
        \(AE_{j,t}=AS_{j,t}/MS_{j,t}\) & Allocative efficiency of asset \(j\)'s continuous double auction when \(MS_{j,t}>0\); reported as zero and excluded from the run mean when no positive-surplus opportunity exists. & Dimensionless ratio \\
        \(k_i\) & Trading rule assigned to chartist \(i\); the same rule is applied to every asset. & Technical rule set \\
        \(L_{j,t}^{\mathrm{BO}},H_{j,t}^{\mathrm{BO}}\) & Previous-window low and high of asset \(j\) used by breakout chartists. & Price units \\
        \(s_{i,j,t}^{C},s_{i,j,t}^{F}\) & Trading signals of chartist and fundamentalist agents. & Buy, sell, or hold states \\
        \(\bar P_{i,j,t}^{\text{entry}}\) & Average execution price of agent \(i\)'s current long or short position in asset \(j\). & Price units \\
        \(\pi_{i,j,t}\) & Unrealised profit-and-loss ratio on asset \(j\). & Dimensionless ratio \\
        \bottomrule
    \end{longtable}

    \subsubsection*{Model Parameters}

    \begin{longtable}{@{}L{0.22\textwidth}L{0.50\textwidth}L{0.20\textwidth}@{}}
        \toprule
        \textbf{Symbol} & \textbf{Meaning} & \textbf{Unit/Domain} \\
        \midrule
        \endfirsthead
        \toprule
        \textbf{Symbol} & \textbf{Meaning} & \textbf{Unit/Domain} \\
        \midrule
        \endhead
        \(r_f\) & Risk-free interest rate per period. & Per-period return \\
        \(\ell_f=\log(1+r_f)\) & Risk-free log return per period; the sure return of the cash alternative in CPT valuation. & Per-period log return \\
        \(\mathcal H_i\) & Forecast horizon of agent \(i\): the number of steps over which its forecast is expected to play out (\(\mathcal H^F\) for fundamentalists, \(\mathcal H^C\) for chartists, set individually for the protagonist). & Real number, \(\mathcal H_i\ge1\) \\
        \(\rho_i\) & Risk-aversion coefficient of agent \(i\). & Positive real number \\
        \(v_{\Sigma}\) & Trailing log-return window used to estimate the shared covariance matrix \(\Sigma_t\). & Positive integer \\
        \(v\) & Lookback window for the legacy single-asset rolling volatility estimate; the implementation allows separate fundamentalist and chartist windows. & Positive integer \\
        \(h_S,h_L\) & Short and long moving-average windows, with \(h_S\le h_L\). & Positive integers \\
        \(h_{\mathrm{RSI}}\) & Lookback window of the RSI indicator; the current implementation uses \(h_{\mathrm{RSI}}=14\). & Positive integer \\
        \(h_{\mathrm{BO}}\) & Lookback window for Donchian-style breakout chartists. & Positive integer \\
        \(\theta^M\) & Fundamentalist mispricing threshold, applied to the log mispricing. & Log-price gap \\
        \(\theta_M\) & Chartist momentum deadband applied to the log moving-average gap. & Log gap, half-open unit interval \\
        \(\theta^{SL}\) & Stop-loss threshold. & Dimensionless loss ratio \\
        \(\theta^{TP}\) & Take-profit threshold. & Dimensionless gain ratio \\
        \(M_0^{S}\) & Initial margin requirement for opening or increasing short exposure. & Positive ratio \\
        \(M_{\min}^{S}\) & Maintenance margin threshold below which a short position is forced to cover. & Positive ratio, \(M_{\min}^{S}\le M_0^{S}\) \\
        \(\varphi^{B}\) & Fraction of an asset's initial share supply that can be borrowed for short selling. & \([0,1]\) \\
        \(f_{\mathrm{tc}}\) & One-way proportional transaction-cost rate paid by each executed trade side. & \([0,1)\) \\
        \(\iota_{\mathrm{tc}}^D\) & Indicator that transaction costs affect discretionary no-trade bands. & \(\{0,1\}\) \\
        \(\iota_{\mathrm{tc}}^B\) & Indicator that transaction costs apply to forced buy-ins. & \(\{0,1\}\) \\
        \(\delta_j\) & Tick size of asset \(j\), used to quantise its limit prices; very small or zero values leave raw quotes unrounded. & Price units, \(\delta_j\ge0\) \\
        \(\tau\) & Hold-tolerance band around the market price. & Log-price gap \\
        \(\kappa\) & Sensitivity of quotes to social opinion. & Non-negative real number \\
        \(\theta_{\mathrm{MO}}\) & Log-price-gap threshold above which a submitted order is converted to market style. & Non-negative log gap \\
        \(\rho_{\mathrm{MO}}\) & Random market-order conversion probability. & \([0,1]\) \\
        \(\theta_{\mathrm{col}}\) & Market-order price collar: maximum absolute log distance from the reference price at which a market order may execute; zero disables the collar. & Non-negative log gap \\
        \(\rho_0,\gamma_{\mathrm{mis}}\) & Base stale-order cancellation probability and sensitivity to quote distance from the latest market price. & \([0,1]\), non-negative sensitivity \\
        \(\varphi_{\mathrm{CPT}}\) & Seeded fraction of agents assigned to use CPT preferences when the global CPT switch is enabled. & \([0,1]\) \\
        \(\lambda_i\) & CPT loss-aversion coefficient of agent \(i\). & Positive real number \\
        \(\alpha_i^G,\alpha_i^L\) & CPT value-function curvature for gains and losses. & Positive real numbers \\
        \(\psi_i^G,\psi_i^L\) & CPT probability-distortion parameters for gains and losses. & Positive real numbers \\
        \(m_{\nu}\) & Configured CPT reference-point mode. & Configured, entry price, peak wealth, moving-average price, or initial wealth \\
        \(\nu_{i,j,t}\) & Behavioural CPT anchor in price or wealth units, recorded for audit. & Positive price or wealth scale \\
        \(\varrho_{i,j,t}\) & Reference log return: the anchor mapped into return space and used by the CPT transform. & \([-0.95,0.95]\) \\
        \(\nu_i^0\) & Configured fallback CPT reference point. & Positive price or wealth scale \\
        \(W_{i,t}^{\max}\) & Highest positive mark-to-market wealth observed by agent \(i\) up to time \(t\). & Currency units \\
        \(h_{\nu}\) & Moving-average window used when \(m_{\nu}\) selects a moving-average price reference. & Positive integer \\
        \(n_{\mathrm{CPT}}\) & Number of scenarios in the moment-matched CPT return ladder; sets only its resolution, not its variance. & Odd integer, \(n_{\mathrm{CPT}}\ge3\) \\
        \(\sigma_N\) & Standard deviation of the noise-trader log-belief shock. & Positive return scale \\
        \(\phi_N\) & Fraction of wealth committed to one noise-trader order. & Positive unit interval \\
        \(\bar w_N^{\text{short}}\) & Maximum absolute short risky weight available to noise traders. & \([0,1]\) \\
        \(\omega_H\) & Total target risky-asset portfolio weight of passive buy-and-hold agents, split equally across assets. & \((0,1]\) \\
        \(\beta_F\) & Legacy fundamental-pull coefficient retained in configuration; inactive under continuous LOB price formation. & Non-negative coefficient \\
        \(\beta_D\) & Legacy order-imbalance impact coefficient retained in configuration; inactive under continuous LOB price formation. & Non-negative coefficient \\
        \bottomrule
    \end{longtable}

    \subsubsection*{Stochastic and Social Variables}

    \begin{longtable}{@{}L{0.22\textwidth}L{0.50\textwidth}L{0.20\textwidth}@{}}
        \toprule
        \textbf{Symbol} & \textbf{Meaning} & \textbf{Unit/Domain} \\
        \midrule
        \endfirsthead
        \toprule
        \textbf{Symbol} & \textbf{Meaning} & \textbf{Unit/Domain} \\
        \midrule
        \endhead
        \(z_{j,t}\) & Independent standard-normal draw for asset \(j\), taken in declared index order. & \(\mathcal{N}(0,1)\) \\
        \(\zeta_{j,t}^{F}\) & Correlated standard-normal fundamental innovation of asset \(j\). & \(\mathcal{N}(0,1)\) \\
        \(\rho_{AB}\) & Pairwise correlation between the standardised fundamental innovations of any two distinct assets. & \([-1,1]\) \\
        \(a,b\) & Coefficients of the one-factor square root of the equicorrelation matrix. & Real numbers \\
        \(\mu_{F,j}\) & Deterministic drift of asset \(j\)'s fundamental process; log-return units under GBM, price units under the legacy arithmetic walk. & Per-period drift \\
        \(\sigma_{F,j}\) & Innovation scale of asset \(j\)'s fundamental process; log-return units under GBM, price units under the legacy arithmetic walk. & Positive scale \\
        \(\varepsilon_{j,t}=\sigma_{F,j}\zeta_{j,t}^{F}\) & Fundamental-value shock of asset \(j\). & Innovation \\
        \(F_{\min,j}\) & Lower bound imposed on asset \(j\)'s latent fundamental value. & Positive price value \\
        \(\beta_{\mathrm{warm}}\) & Fundamental-pull coefficient used only in the warm-up price process. & Non-negative coefficient \\
        \(\sigma_{\mathrm{warm}}\) & Standard deviation of the warm-up market-price noise. & Per-period return \\
        \(\xi_u^{\mathrm{warm}}\) & Warm-up price-process shock. & Return innovation \\
        \(\xi_t\) & Legacy exogenous market-price shock; set to zero under continuous LOB price formation. & Return innovation \\
        \(\sigma_P\) & Legacy market-price shock scale retained in configuration; inactive under continuous LOB price formation. & Per-period return \\
        \(\zeta_{i,t}\) & Noise-trader log-belief shock, \(\zeta_{i,t}\sim\mathcal{N}(0,\sigma_N^2)\). & Log-price shock \\
        \(\zeta_{\max}\) & Numerical bound applied to noise-trader log quote shocks before exponentiation. & Positive real number \\
        \(U_{i,t}\) & Uniform draw selecting which asset a noise trader trades in a multi-asset market. & \(\mathcal{U}(0,1)\) \\
        \((u_i,v_i)\) & Spatial coordinates of agent \(i\). & Synthetic social space \\
        \(d_{il}\) & Euclidean distance between agents \(i\) and \(l\). & Distance units \\
        \(r_{\mathrm{soc}}\) & Social-interaction radius. & Distance units \\
        \(\mathcal{N}_i\) & Neighbourhood of agent \(i\) within radius \(r_{\mathrm{soc}}\). & Set of neighbouring agents \\
        \(\omega_{il}\) & Distance-based social weight assigned by agent \(i\) to neighbour \(l\). & Positive weight \\
        \(\eta_{i,t}\) & Private signal of agent \(i\), \(\eta_{i,t}\sim\mathcal{N}(0,\sigma_{i,\text{priv}}^2)\). & Return signal \\
        \(\sigma_{i,\text{priv}}\) & Standard deviation of agent \(i\)'s private signal. & Return standard deviation \\
        \(S_{i,t}\) & Weighted-average social signal observed by agent \(i\). & Return signal \\
        \(o_{i,t}\) & Opinion of agent \(i\) after combining private and social signals; one opinion per agent, applied to every asset. & Return signal \\
        \(\gamma_i\) & Social-influence weight of agent \(i\). & \([0,1]\) \\
        \(\tilde P_{i,j,t}\) & Pre-quantisation valuation quote of agent \(i\) for asset \(j\). & Price units \\
        \(P_{i,j,t}^{\text{lim}}\) & Effective submitted limit price after tick-size quantisation. & Price units \\
        \bottomrule
    \end{longtable}

    \subsubsection*{Market Microstructure}

    \begin{longtable}{@{}L{0.22\textwidth}L{0.50\textwidth}L{0.20\textwidth}@{}}
        \toprule
        \textbf{Symbol} & \textbf{Meaning} & \textbf{Unit/Domain} \\
        \midrule
        \endfirsthead
        \toprule
        \textbf{Symbol} & \textbf{Meaning} & \textbf{Unit/Domain} \\
        \midrule
        \endhead
        \(P_{i,j,t}^{\text{lim}}\) & Effective limit price submitted by agent \(i\) in asset \(j\). & Price units \\
        \(\sigma_{i,t}^{\text{ord}}\) & Order style. & \(\{\mathrm{Limit},\mathrm{Market}\}\) \\
        \(\mathrm{asset}(\mathfrak{o})\) & Asset index carried by an order; routes it to exactly one book. & Asset identifier \\
        \(\mathrm{owner}(\mathfrak{o})\) & Agent identifier of the order's owner, used for self-trade prevention. & Agent identifier \\
        \(q_{i,t}^{\text{rem}}\) & Remaining unfilled order quantity. & Shares or units \\
        \(T_{i,t}^{0},T_{i,t}^{X}\) & Order creation and expiry time steps. & Non-negative integers \\
        \(n_{i,t}^{\text{seq}}\) & Stable sequence number used for FIFO priority. & Positive integer \\
        \(\mathcal{L}_{j,t}^B,\mathcal{L}_{j,t}^A\) & Persistent bid and ask books of asset \(j\) before new submissions in period \(t\). & Ordered order sets \\
        \(P_{j,t}^B,P_{j,t}^A\) & Best bid and best ask of asset \(j\) after period-\(t\) processing. & Price units or zero if absent \\
        \(S_{j,t}^{BA}=P_{j,t}^A-P_{j,t}^B\) & Bid-ask spread when both best quotes exist. & Price units \\
        \(M_{j,t}^{Q}=(P_{j,t}^B+P_{j,t}^A)/2\) & Midquote when both best quotes exist. & Price units \\
        \(D_{j,t}^B(h),D_{j,t}^A(h)\) & Visible bid and ask depth across the first \(h\) price levels. & Shares or units \\
        \(Q_{j,t}^{\text{exec}}\) & Volume executed endogenously in asset \(j\)'s limit order book. & Shares or units \\
        \(\widetilde Q_{j,t}\) & Reported traded volume after adding maintenance-margin forced buy-ins. & Shares or units \\
        \(G_k=m_kp_k\) & Gross notional value of execution \(k\). & Currency units \\
        \(F_k^{\mathrm{tc}}\) & One-side transaction cost charged on execution \(k\). & Currency units \\
        \(\chi_{j,t}\) & Indicator that at least one trade is executed in asset \(j\) in period \(t\). & \(\{0,1\}\) \\
        \(P_{j,t}^{\text{last}}\) & Last endogenous transaction price in asset \(j\) in period \(t\). & Price units \\
        \(e_{i,j,t}\) & Executed quantity of agent \(i\) in asset \(j\) at time \(t\). & Shares or units \\
        \(\mathrm{ND}_t\) & Submitted net demand before execution. & Shares or units \\
        \(\mathrm{NI}_t\) & Submitted notional imbalance before execution. & Currency units \\
        \(\bar I_t\) & Normalised notional imbalance. & \([-1,1]\) \\
        \(g_t^F\) & Relative fundamental-value signal recorded as a diagnostic. & Dimensionless ratio \\
        \(r_{j,t}^{\mathrm{raw}}\) & Realised one-period simple return before storing \(P_{j,t+1}\); zero when no trade is executed. & \([-0.95,0.95]\) \\
        \(\ell_{j,t+1}\) & Realised one-period log return of asset \(j\). & Log return \\
        \(\epsilon_P\) & Minimum price scale used for numerical protection. & Small positive price \\
        \(\epsilon_q\) & Minimum absolute fill quantity treated as a real execution. & Small positive quantity \\
        \(\varphi_{\mathrm{fill}}\) & Minimum economic fill expressed as a fraction of the requested match quantity. & \([0,1)\) \\
        \(m^{\mathrm{req}},m^{\min}\) & Requested match quantity and minimum economic fill of a candidate execution. & Shares or units \\
        \(K_{\mathrm{match}}\) & Terminal cap on matching-loop iterations for one incoming order. & Positive integer \\
        \bottomrule
    \end{longtable}

    \section{Simulation Algorithm}\label{sec:simulation_algorithm}

    The following pseudocode summarises the simulation loop implemented by the
    model. Its purpose is to make the execution order transparent before the
    individual equations are introduced.

    \begin{framed}
                \raggedright
                \textbf{Inputs:} simulation configuration, deterministic seed, the enabled
                asset set \(\mathcal{J}\) with per-asset initial prices \(P_{j,0}\) and
                fundamental values \(F_{j,0}=P_{j,0}\), the cross-asset correlation
                \(\rho_{AB}\), agent population, and output paths.

                \textbf{Initialisation}
                \begin{enumerate}
                    \item Normalise the configuration: drop disabled assets from
                    \(\mathcal{J}\) into \(\mathcal{J}^{0}\), fix the declared index order,
                    and resolve the primary asset index \(j^{\ast}\).
                    \item Initialise each asset's market and fundamental histories with
                    \(P_{j,0}\) and \(F_{j,0}\).
                    \item Initialise fundamentalist, chartist, noise-trader, and passive
                    buy-and-hold agents with cash, per-asset inventory, strategy parameters, CPT
                    preferences, and personal reference-point state where applicable.
                    \item If social interaction is enabled, assign spatial coordinates,
                    private-signal noise scales, social-influence weights, and neighbourhoods.
                    \item Generate warm-up observations with the configured temporary price
                    process, advancing all fundamentals jointly so their innovations carry the
                    configured correlation, so technical indicators and covariance estimates
                    are well defined before the analysed continuous-LOB phase.
                \end{enumerate}

                \textbf{For each simulation step \(t=1,\ldots,T\)}
                \begin{enumerate}
                    \item Accrue risk-free interest on agent cash holdings.
                    \item Advance every asset's latent fundamental value \(F_{j,t}\) in one
                    joint draw with correlated innovations.
                    \item Update private and social opinions \(o_{i,t}\) for all agents.
                    \item Expire old resting orders and cancel stale resting orders from each
                    asset's persistent bid and ask books.
                    \item Estimate the shared log-return covariance matrix \(\Sigma_t\) over
                    the trailing window \(v_{\Sigma}\).
                    \item Select active agents and collect orders. Portfolio-optimising agents
                    solve one joint Markowitz problem over all assets and emit at most one
                    order per asset from their strategy valuation, CPT-adjusted return when
                    applicable, and social quote overlay when enabled; noise traders draw one
                    asset uniformly. Sufficiently aggressive valuation gaps, or a
                    configured random event, convert an order from limit style to market style.
                    \item Match incoming orders continuously against the persistent
                    price-time-priority limit order book \emph{of their own asset}; books never
                    cross assets. A resting order owned by the incoming order's own agent is
                    cancelled rather than traded with. Market orders whose execution price
                    breaches the collar are rejected. Candidate fills below the
                    minimum economic fill threshold are rejected and the binding side
                    is cancelled. Executions update cash net of
                    proportional transaction costs, per-asset inventory, short-sale state, trade
                    counts, per-asset average entry-price references, cumulative
                    transaction-cost diagnostics, and the trade ledger immediately.
                    \item For each asset independently: if at least one endogenous trade is
                    executed, store the last transaction price as \(P_{j,t+1}\); otherwise carry
                    forward the previous observed market price.
                    \item Refresh agent wealth and personal peak wealth across all assets,
                    enforce short-sale maintenance margin per asset through forced buy-ins when
                    required, refresh wealth again when a forced buy-in has traded, and record
                    step-level outputs: the primary asset's single-asset row plus one
                    per-\((\text{step},\text{asset})\) row holding price, fundamental value,
                    returns, book state, cancellation counters, allocative efficiency, mean
                    agent portfolio weight, traded volume, and trade-ledger entries.
                \end{enumerate}

                \textbf{After the final step:} export simulation details, agent snapshots,
                CPT/reference-point and transaction-cost diagnostics, trade-ledger records
                when enabled, summary metrics, per-asset stylised-fact diagnostics, the
                long-format per-asset time series, resolved configuration, metadata, and the
                experiment manifest.
    \end{framed}

    \section{Fundamental Value Generation}\label{sec:fundamental_value}

    \subsection{Per-Asset Process}

    Every asset's latent fundamental value follows geometric Brownian motion in
    Euler-Maruyama log form. This is the model's fundamental-value theory and is
    the default for every declared asset:
    \[
        F_{j,t+1}
        =
        \max\!\left\{
            F_{\min,j},\,
            F_{j,t}\exp\!\left(
                \left(\mu_{F,j}-\tfrac{1}{2}\sigma_{F,j}^{2}\right)
                +\sigma_{F,j}\zeta_{j,t}^{F}
            \right)
        \right\},
        \qquad
        \zeta_{j,t}^{F}\sim\mathcal{N}(0,1).
    \]
    The It\^o correction \(-\tfrac12\sigma_{F,j}^2\) keeps \(\mu_{F,j}\)
    interpretable as the expected arithmetic growth rate, so
    \(\mathbb{E}[F_{j,t+1}\mid F_{j,t}]=F_{j,t}e^{\mu_{F,j}}\) while
    \(\mathbb{E}[\log F_{j,t+1}-\log F_{j,t}]=\mu_{F,j}-\tfrac12\sigma_{F,j}^2\).
    This follows the standard Black-Scholes discretisation.

    Two properties matter for the multi-asset design. First, multiplicative
    dynamics keep the value strictly positive without relying on the floor
    \(F_{\min,j}>0\), which now guards only against extreme downward excursions
    reaching the price grid. Second, \(\mu_{F,j}\) and \(\sigma_{F,j}\) are in
    per-step \emph{log-return} units, so they are directly comparable across
    assets regardless of price level: an asset trading at \(100\) and one trading
    at \(40\) with the same \(\sigma_{F,j}\) carry the same proportional risk.
    Under the previous additive specification the same numeric shock meant a very
    different economic disturbance at different price levels, which would make
    cross-asset comparison meaningless.

    A legacy arithmetic random walk,
    \[
        F_{j,t+1}=\max\{F_{\min,j},\,F_{j,t}+\mu_{F,j}+\varepsilon_{j,t}\},
        \qquad
        \varepsilon_{j,t}=\sigma_{F,j}\zeta_{j,t}^{F},
    \]
    with drift and shock in \emph{price} units, is retained and must be requested
    explicitly per asset. It exists only so that studies predating the geometric
    process can be reproduced exactly.

    \subsection{Cross-Asset Correlation}\label{sec:fundamental_correlation}

    Fundamentals in a multi-asset market are not independent: a common
    macroeconomic factor makes them co-move, and that co-movement is precisely
    what limits the diversification available to agents. All assets are therefore
    advanced in one joint draw. Let \(R\) be the equicorrelation matrix with unit
    diagonal and off-diagonal \(\rho_{AB}\in[-1,1]\). One standard normal
    \(z_{j,t}\) is drawn per asset in declared index order and mapped to
    correlated standard normals by the closed-form one-factor square root of
    \(R\):
    \[
        \zeta_{j,t}^{F}
        =
        a\,z_{j,t}+b\sum_{k\in\mathcal{J}}z_{k,t},
        \qquad
        a=\sqrt{1-\rho_{AB}},
        \qquad
        b=\frac{\sqrt{1+(n-1)\rho_{AB}}-a}{n}.
    \]
    Then \(\Corr(\zeta_{j,t}^{F},\zeta_{k,t}^{F})=\rho_{AB}\) for \(j\ne k\), and
    each asset applies its own process to its own \(\zeta_{j,t}^{F}\). The
    closed form is preferred to a generic Cholesky factorisation because it stays
    exact on the positive-semi-definite boundary --- at \(\rho_{AB}=1\), or at the
    lower admissible limit \(\rho_{AB}=-1/(n-1)\) --- where Cholesky would fail on
    a singular matrix. Correlations outside \([-1,1]\), or values implying a
    non-positive-semi-definite \(R\), are rejected at configuration time.

    With \(n=1\), or with \(\rho_{AB}=0\), the map is exactly the identity and is
    short-circuited, so the innovation \(\sigma_{F,j}z_{j,t}\) is bit-identical to
    an independent single-asset draw and the historical random-draw sequence is
    preserved.

    In the two-asset case \(\rho_{AB}\) is exactly \(\rho_{AB}\) between assets
    \(A\) and \(B\); as \(\rho_{AB}\to1\) the Markowitz rule of
    Section~\ref{sec:decision_rules} can no longer reduce risk by spreading
    across assets.

    \subsection{Warm-Up Price Process}

    During warm-up only, the implementation generates a temporary market-price
    history for each asset so the first analysed order-book step has enough
    observations for moving averages, rolling volatility, and the covariance
    matrix. Fundamentals still advance jointly and correlated; the warm-up price
    noise is drawn from a separate engine. Let \(r_{j,u}^{\mathrm{warm}}\) denote
    the warm-up simple return of asset \(j\). The warm-up price process is
    \[
        r_{j,u}^{\mathrm{warm}}
        =
        \clip\!\left(
                   \beta_{\mathrm{warm}}\frac{F_{j,u}-P_{j,u}}{\max\{P_{j,u},\epsilon_P\}}
                   +\xi_{j,u}^{\mathrm{warm}};
            -0.95,0.95
        \right),
        \qquad
        \xi_{j,u}^{\mathrm{warm}}\sim\mathcal{N}(0,\sigma_{\mathrm{warm}}^2),
    \]
    \[
        P_{j,u+1}=\max\{\epsilon_P,P_{j,u}(1+r_{j,u}^{\mathrm{warm}})\}.
    \]
    These warm-up observations are not interpreted as analysed LOB transactions;
    they only initialise state variables required by the trading rules.

    \section{Agent Decision Rules}\label{sec:decision_rules}

    At each time step \(t\), agent \(i\) is active with probability
    \[
        \Pr(a_{i,t}=1)=\alpha.
    \]
    If inactive, the agent submits no order.

    Let \(P_{j,t}\) be the latest market price of asset \(j\), \(F_{j,t}\) its
    latest fundamental price, \(c_{i,t}\) the agent's single shared cash account,
    \(x_{i,j,t}\) its holding of asset \(j\), and
    \[
        W_{i,t}=c_{i,t}+\sum_{j\in\mathcal{J}}x_{i,j,t}P_{j,t}
    \]
    the current wealth. Wealth is a portfolio quantity: an agent funds purchases in
    any asset from the same cash balance and is marked to market against every
    position it holds. Define
    \[
        \clip(z;\underline z,\overline z)=\min\{\overline z,\max\{\underline z,z\}\},
        \qquad
        \pospart{z}=\max\{z,0\}.
    \]

    Proportional transaction costs are represented by the one-way rate
    \(f_{\mathrm{tc}}\). The cash required to fund one purchased share at price
    \(p\) is \(p(1+f_{\mathrm{tc}})\), so the cost-aware affordable buy quantity is
    \[
        A_{i,t}^{B}(p)
        =
        \frac{c_{i,t}}{p(1+f_{\mathrm{tc}})}.
    \]
    When transaction costs affect decisions, the strategy-level no-trade band is
    enlarged by the round-trip hurdle
    \[
        \tau_{\mathrm{tc}}^{\log}
        =
        \log\frac{1+f_{\mathrm{tc}}}{1-f_{\mathrm{tc}}},
        \qquad
        \tau_{\mathrm{tc}}^{r}
        =
        \frac{2f_{\mathrm{tc}}}{1-f_{\mathrm{tc}}}.
    \]
    The log hurdle is used for valuation-price gaps, while the simple-return
    hurdle is used for the fundamentalist mispricing signal. If the
    decision-cost switch is disabled, both hurdles are treated as zero in the
    trading rules.

    \subsection{Risk Estimation}\label{sec:covariance}

    All agents share one estimate of market risk, computed once per step from
    trailing one-period log returns of the stored market prices. Warm-up history
    is included, so the estimate exists from the first trading step. Let
    \[
        \ell_{j,u}=\log\frac{P_{j,u}}{P_{j,u-1}}
    \]
    and let \(T_t=\min\{v_{\Sigma},\,\text{shortest available return history}\}\)
    be the common trailing sample length, which keeps the matrix internally
    consistent across assets. With \(\bar\ell_j\) the sample mean over that
    window, the entries of the covariance matrix are
    \[
        \Sigma_{t,jk}
        =
        \frac{1}{T_t-1}
        \sum_{u=t-T_t+1}^{t}
        \left(\ell_{j,u}-\bar\ell_j\right)\left(\ell_{k,u}-\bar\ell_k\right),
        \qquad
        \sigma_{j,t}=\sqrt{\Sigma_{t,jj}}.
    \]
    Fewer than two common observations give the zero matrix; callers regularise
    before factorisation in either case. Using log-return covariance keeps the
    allocation rule consistent with the return series used in the stylised-fact
    diagnostics, and makes the risk estimate comparable across assets trading at
    different price levels.

    \subsection{Portfolio Allocation}\label{sec:allocation}

    Each strategy forecast \(\breve r_{i,j,t}\) is the excess log return the agent
    expects over its forecast horizon \(\mathcal H_i\ge1\), the number of steps
    over which it expects the forecast to play out: a mispricing, for instance,
    is not expected to close in a single step. The Markowitz rule needs the
    expected return and the covariance over the same period, and the covariance
    is per step (in Merton's rule both are rates per unit time), so the forecast
    is spread over the horizon. Given the vector of prospect-adjusted excess
    returns
    \(\breve{\mathbf r}_{i,t}=(\breve r_{i,1,t},\ldots,\breve r_{i,n,t})^{\top}\),
    the agent solves the full Markowitz problem
    \[
        \mu_{i,j,t}=r_f+\frac{\breve r_{i,j,t}}{\mathcal H_i},
        \qquad
        \mathbf w_{i,t}^{\mathrm{raw}}
        =
        \frac{1}{\rho_i\mathcal H_i}\Sigma_t^{-1}\breve{\mathbf r}_{i,t},
    \]
    where \(\rho_i>0\) is the agent's risk aversion. Equivalently, the agent
    weighs the horizon return \(\breve{\mathbf r}_{i,t}\) against the horizon risk
    \(\mathcal H_i\Sigma_t\) of independent per-step returns. Scaling the
    fundamentalist forecast by a mean-reversion time in this way follows
    Chiarella, Iori and Perell\'o's order-book model. With \(\mathcal H_i=1\) a
    whole mispricing is treated as a one-step return; at the model's typical
    volatility nearly every raw weight then falls outside its bounds, so risk
    aversion and the covariance stop shaping positions. Setting
    \(\mathcal H_i=1\) reproduces that one-step rule exactly. The residual weight
    \(1-\sum_j w_{i,j,t}\) is held at the risk-free rate implicitly, as in the
    single-asset model. This is an intentionally transparent mean-variance
    allocation rule in the tradition of Markowitz; the behavioural contribution
    enters through the forecast and preference transforms rather than through a
    separate utility-maximisation problem. Cross-asset covariance now shapes the
    demand for each asset: a forecast on one asset moves the position in the
    other through \(\Sigma_t^{-1}\), and the diversification benefit shrinks as
    \(\rho_{AB}\) rises.

    The solve is implemented by Cholesky factorisation of a regularised copy of
    \(\Sigma_t\): a variance floor is applied to the diagonal, and if
    factorisation fails an escalating diagonal loading (starting at \(10^{-6}\)
    and multiplied by ten, up to three retries) is added. If factorisation still
    fails, or produces a non-finite solution, the rule degrades gracefully to the
    diagonal solve
    \[
        w_{i,j,t}^{\mathrm{raw}}
        =
        \frac{\breve r_{i,j,t}}{\rho_i\mathcal H_i\sigma_{j,t}^{2}},
    \]
    which for \(n=1\) is exactly the historical scalar mean-variance rule.

    Constraints are then applied as a projection. Weights are clamped to the
    per-asset interval first; if gross exposure then exceeds the cap, all weights
    are scaled proportionally and clamped once more:
    \[
        w_{i,j,t}^{(1)}=\clip\!\left(w_{i,j,t}^{\mathrm{raw}};\underline w_{i,j},\overline w_{i,j}\right),
        \qquad
        \Lambda_i=\sum_{j}\left|w_{i,j,t}^{(1)}\right|,
    \]
    \[
        w_{i,j,t}=
        \begin{cases}
            w_{i,j,t}^{(1)}, & \Lambda_i\le G_i \ \text{or}\ G_i\le 0,\\[1.2ex]
            \clip\!\left(\dfrac{G_i}{\Lambda_i}w_{i,j,t}^{(1)};\underline w_{i,j},\overline w_{i,j}\right), & \Lambda_i>G_i.
        \end{cases}
    \]
    The gross-exposure cap \(G_i\) is what prevents a two-asset portfolio from
    silently doubling an agent's market exposure relative to the single-asset
    baseline. The target position in each asset is
    \[
        x_{i,j,t}^{\ast}=\frac{W_{i,t}w_{i,j,t}}{P_{j,t}}.
    \]
    For agents not assigned to Cumulative Prospect Theory (CPT),
    \(\breve r_{i,j,t}=\hat r_{i,j,t}\). For agents assigned to CPT,
    \(\breve r_{i,j,t}\) is the prospect-adjusted excess return defined in
    Section~\ref{sec:cpt}.

    Because all buys are funded from one cash account, order generation carries a
    running budget \(\tilde c_{i,t}\), initialised at \(c_{i,t}\) and decremented
    by \(P_{i,j,t}^{\text{lim}}q_{i,j,t}(1+f_{\mathrm{tc}})\) as each buy order is
    queued. This stops simultaneous buys across assets from jointly
    overcommitting the same cash at submission time; the clearing engine still
    enforces the hard affordability constraint at execution time
    (Section~\ref{sec:short_sale_constraints}).

    \subsection{Technical Indicators}

    For chartists, define the moving average of asset \(j\) with horizon
    \(h\in\{h_S,h_L\}\) by
    \[
        \MA_{j,t}^{(h)}=\frac{1}{h}\sum_{u=t-h+1}^{t}P_{j,u}.
    \]
    Let \(\theta_M\) denote the configured non-negative momentum-signal threshold,
    strictly below one. It is
    implemented as a deadband around zero so very small moving-average gaps are not
    interpreted as economically meaningful trend-following evidence. The gap is
    measured in log space,
    \[
        g_{j,t}^M=\log\frac{\MA_{j,t}^{(h_S)}}{\MA_{j,t}^{(h_L)}},
    \]
    so that it is commensurate with the log-return volatility used by the
    allocation rule. Both the gap and the ratios below are defined as zero
    whenever either side falls below the minimum price scale \(\epsilon_P\); this
    is a numerical guard, not an economic assumption.

    For RSI, let \(\Delta P_{j,u}=P_{j,u}-P_{j,u-1}\), and
    \[
        \bar G_{j,t}=\frac{1}{h_{\mathrm{RSI}}}\sum_{u=t-h_{\mathrm{RSI}}+1}^{t}\pospart{\Delta P_{j,u}},
        \qquad
        \bar L_{j,t}=\frac{1}{h_{\mathrm{RSI}}}\sum_{u=t-h_{\mathrm{RSI}}+1}^{t}\pospart{-\Delta P_{j,u}},
    \]
    \[
        \RSI_{j,t}=
        \begin{cases}
            100, & \bar L_{j,t}=0,\\
            0, & \bar G_{j,t}=0,\\
            100-\dfrac{100}{1+\bar G_{j,t}/\bar L_{j,t}}, & \text{otherwise}.
        \end{cases}
    \]
    For Donchian-style breakout chartists, the latest price is compared with the
    channel formed by the previous \(h_{\mathrm{BO}}\) observations, excluding the
    current price:
    \[
        L_{j,t}^{\mathrm{BO}}=\min\{P_{j,t-h_{\mathrm{BO}}},\ldots,P_{j,t-1}\},
        \qquad
        H_{j,t}^{\mathrm{BO}}=\max\{P_{j,t-h_{\mathrm{BO}}},\ldots,P_{j,t-1}\}.
    \]
    Excluding \(P_{j,t}\) prevents the current observation from defining its own
    breakout threshold.

    \subsection{Chartists}
    Each chartist \(i\) is assigned one strategy
    \(k_i\in\{\mathrm{MOM},\mathrm{MR},\mathrm{RSI},\mathrm{BO}\}\) and applies that
    same rule to every asset. Its technical state in asset \(j\) is
    summarised by the discrete signal
    \[
        s_{i,j,t}^{C}=
        \begin{cases}
            \mathrm{Buy}, & k_i=\mathrm{MOM}\ \text{and}\ g_{j,t}^M>\theta_M,\\
            \mathrm{Sell}, & k_i=\mathrm{MOM}\ \text{and}\ g_{j,t}^M<-\theta_M,\\
            \mathrm{Buy}, & k_i=\mathrm{MR}\ \text{and}\ P_{j,t}<\MA_{j,t}^{(h_L)},\\
            \mathrm{Sell}, & k_i=\mathrm{MR}\ \text{and}\ P_{j,t}>\MA_{j,t}^{(h_L)},\\
            \mathrm{Buy}, & k_i=\mathrm{RSI}\ \text{and}\ \RSI_{j,t}<30,\\
            \mathrm{Sell}, & k_i=\mathrm{RSI}\ \text{and}\ \RSI_{j,t}>70,\\
            \mathrm{Buy}, & k_i=\mathrm{BO}\ \text{and}\ P_{j,t}>H_{j,t}^{\mathrm{BO}},\\
            \mathrm{Sell}, & k_i=\mathrm{BO}\ \text{and}\ P_{j,t}<L_{j,t}^{\mathrm{BO}},\\
            \mathrm{Hold}, & \text{otherwise}.
        \end{cases}
    \]

    This signal records the direction implied by the technical rule. In the current
    implementation, actual order side is determined by the valuation quote in
    Section~\ref{sec:valuation_quotes}; therefore a chartist can submit no order
    even after a non-hold technical signal if the valuation quote lies inside the
    hold-tolerance band.

    Its expected excess return on asset \(j\) over the chartist forecast horizon
    \(\mathcal H^C\) is
    \[
        \hat r_{i,j,t}^{C}=
        \begin{cases}
            0, & k_i=\mathrm{MOM}\ \text{and}\ |g_{j,t}^M|\le\theta_M,\\[1ex]
            \clip\!\left(g_{j,t}^M;-1,1\right), & k_i=\mathrm{MOM}\ \text{and}\ |g_{j,t}^M|>\theta_M,\\[1ex]
            \clip\!\left(\log\dfrac{\MA_{j,t}^{(h_L)}}{P_{j,t}};-1,1\right), & k_i=\mathrm{MR},\\[1ex]
            \clip\!\left(0.05\,\dfrac{50-\RSI_{j,t}}{50};-1,1\right), & k_i=\mathrm{RSI},\\[1ex]
            \clip\!\left(\log\dfrac{P_{j,t}}{H_{j,t}^{\mathrm{BO}}};-1,1\right), & k_i=\mathrm{BO}\ \text{and}\ P_{j,t}>H_{j,t}^{\mathrm{BO}},\\[1ex]
            \clip\!\left(\log\dfrac{P_{j,t}}{L_{j,t}^{\mathrm{BO}}};-1,1\right), & k_i=\mathrm{BO}\ \text{and}\ P_{j,t}<L_{j,t}^{\mathrm{BO}},\\[1ex]
            0, & k_i=\mathrm{BO}\ \text{and}\ L_{j,t}^{\mathrm{BO}}\le P_{j,t}\le H_{j,t}^{\mathrm{BO}}.
        \end{cases}
    \]
    Every branch except RSI is a log ratio, so the forecast is a log return and is
    directly comparable with the entries of \(\Sigma_t\). The RSI branch is already
    a bounded dimensionless score and is left unchanged.
    The excess return used in allocation is
    \begin{align*}
        \breve r_{i,j,t}^{C}&=
        \begin{cases}
            \hat r_{i,j,t}^{C}, & \text{agent } i \text{ is not assigned CPT},\\
            \mathcal{T}_i\!\left(\hat r_{i,j,t}^{C},\sigma_{j,t};\,\varrho_{i,j,t}\right), & \text{agent } i \text{ is assigned CPT},
        \end{cases}
        \\
        \mu_{i,j,t}^{C}&=r_f+\breve r_{i,j,t}^{C}/\mathcal H^{C}.
    \end{align*}

    The chartist target weight in every asset is bounded by
    \[
        \underline w_{i,j}=\max\left\{\underline w_i^{\text{cfg}},\,-\bar w_i^{\text{short}}\right\},
        \qquad
        \overline w_{i,j}=\pospart{\overline w_i^{\text{cfg}}},
    \]
    so the configured allocation floor and the short-exposure cap bind
    independently, and the tighter of the two applies.

    Given an order side from the valuation quote, the chartist order quantity in
    asset \(j\) is
    \[
        q_{i,j,t}^{C}=
        \begin{cases}
            \min\left\{\pospart{x_{i,j,t}^{\ast}-x_{i,j,t}},\tilde A_{i,t}^{B}(P_{i,j,t}^{\text{lim}})\right\}, & \mathrm{side}_{i,j,t}=\mathrm{Buy},\\[1.2ex]
            \pospart{x_{i,j,t}-x_{i,j,t}^{\ast}}, & \mathrm{side}_{i,j,t}=\mathrm{Sell},\\
            0, & \text{quote lies inside the hold-tolerance band},
        \end{cases}
    \]
    where \(\tilde A_{i,t}^{B}(p)=\tilde c_{i,t}/\bigl(p(1+f_{\mathrm{tc}})\bigr)\)
    is the affordable buy quantity against the running cash budget of
    Section~\ref{sec:allocation}.

    \subsection{Fundamentalists}
    Define the log mispricing of asset \(j\),
    \[
        m_{j,t}=\log\frac{F_{j,t}}{P_{j,t}}.
    \]
    Positive \(m_{j,t}\) means the latent fundamental value lies above the market
    price, so the asset is treated as undervalued. The log form replaces the
    former relative difference \((F_{j,t}-P_{j,t})/P_{j,t}\) and keeps the
    forecast commensurate with the log-return covariance used for sizing.

    Let \(\bar P_{i,j,t}^{\text{entry}}\) denote the average entry price of agent
    \(i\)'s current long position in asset \(j\) when one exists. The unrealised
    P\&L ratio used for stop-loss and take-profit rules is evaluated per asset
    against that asset's own entry price:
    \[
        \pi_{i,j,t}=\frac{P_{j,t}-\bar P_{i,j,t}^{\text{entry}}}{\bar P_{i,j,t}^{\text{entry}}}.
    \]

    The fundamentalist signal in asset \(j\) is
    \[
        s_{i,j,t}^{F}=
        \begin{cases}
            (\mathrm{Sell},\mathrm{Force}), & x_{i,j,t}>0\ \text{and}\ \pi_{i,j,t}\le -\theta^{SL},\\
        (\mathrm{Sell},\mathrm{Force}), & x_{i,j,t}>0\ \text{and}\ \pi_{i,j,t}\ge \theta^{TP},\\
        (\mathrm{Buy},\mathrm{Rebalance}), & m_{j,t}\ge \theta^{M}+\iota_{\mathrm{tc}}^D\tau_{\mathrm{tc}}^{r},\\
        (\mathrm{Sell},\mathrm{Rebalance}), & m_{j,t}\le -\left(\theta^{M}+\iota_{\mathrm{tc}}^D\tau_{\mathrm{tc}}^{r}\right),\\
        (\mathrm{Hold},\mathrm{None}), & |m_{j,t}|<\theta^{M}+\iota_{\mathrm{tc}}^D\tau_{\mathrm{tc}}^{r}.
        \end{cases}
    \]
    Forced liquidation is therefore asset-specific: a stop-loss in one asset does
    not liquidate the agent's position in the other.

    The forced-liquidation part of this signal is active in the current
    valuation-quote implementation. The mispricing threshold remains part of the
    generated strategy signal, but non-forced buy, sell, and hold decisions in the
    active order path are governed by the valuation quote and hold-tolerance
    condition in Section~\ref{sec:valuation_quotes}.

    Its expected excess return over the fundamentalist forecast horizon
    \(\mathcal H^F\), the number of steps over which the mispricing is expected
    to close, is
    \begin{align*}
        \hat r_{i,j,t}^{F}&=\clip(m_{j,t};-1,1),\\
        \breve r_{i,j,t}^{F}&=
        \begin{cases}
            \hat r_{i,j,t}^{F}, & \text{agent } i \text{ is not assigned CPT},\\
            \mathcal{T}_i\!\left(\hat r_{i,j,t}^{F},\sigma_{j,t};\,\varrho_{i,j,t}\right), & \text{agent } i \text{ is assigned CPT},
        \end{cases}
        \\
        \mu_{i,j,t}^{F}&=r_f+\breve r_{i,j,t}^{F}/\mathcal H^{F}.
    \end{align*}
    The reference log return \(\varrho_{i,j,t}\) is resolved by the configured
    agent-specific reference-point rule in Section~\ref{sec:cpt_reference_points}.

    Fundamentalists remain long-only, so their weight bounds are floored at zero,
    \[
        \underline w_{i,j}=\pospart{\underline w_i^{\text{cfg}}},
        \qquad
        \overline w_{i,j}=\pospart{\overline w_i^{\text{cfg}}},
    \]
    and the desired position is
    \[
        x_{i,j,t}^{\ast,F}=\max\left\{0,\frac{W_{i,t}w_{i,j,t}}{P_{j,t}}\right\},
    \]
    with \(x_{i,j,t}^{\ast,F}=0\) under forced liquidation. Short exposure is
    handled by chartist and noise-trader rules plus the global risk-control layer.

    Given an order side from the valuation quote, the fundamentalist order quantity
    in asset \(j\) is
    \[
        q_{i,j,t}^{F}=
        \begin{cases}
            \min\left\{\pospart{x_{i,j,t}^{\ast,F}-x_{i,j,t}},\tilde A_{i,t}^{B}(P_{i,j,t}^{\text{lim}})\right\},
            & \mathrm{side}_{i,j,t}=\mathrm{Buy},\\[1.2ex]
            x_{i,j,t},
            & s_{i,j,t}^{F}=(\mathrm{Sell},\mathrm{Force}),\\[1.2ex]
            \pospart{x_{i,j,t}-x_{i,j,t}^{\ast,F}},
            & \mathrm{side}_{i,j,t}=\mathrm{Sell}\ \text{and not forced},\\[1.2ex]
            0,
            & \text{quote lies inside the hold-tolerance band}.
        \end{cases}
    \]

    \subsection{Noise Traders}

    Noise traders \(i\in\mathcal{A}^N\) represent uninformed random-liquidity
    demand. They do not optimise a portfolio, so in a multi-asset market each
    active noise trader first selects a single asset uniformly at random from its
    own private random stream,
    \[
        j_{i,t}
        =
        \min\left\{\left\lfloor n\,U_{i,t}\right\rfloor,\,n-1\right\},
        \qquad
        U_{i,t}\sim\mathcal{U}(0,1),
    \]
    and trades only that asset in that step. The draw is taken only when
    \(n>1\), so single-asset private random streams are unchanged. Below,
    \(j\) denotes the selected asset.

    At each active decision time they draw an idiosyncratic log-belief shock
    \[
        \zeta_{i,t}\sim\mathcal{N}(0,\sigma_N^2).
    \]
    The shock is interpreted as a short-run erroneous belief about the market price,
    consistent with noise-trader-risk models in which stochastic beliefs can move
    prices and create limits to arbitrage. For numerical protection, the
    implementation clips the log quote shock before exponentiation:
    \[
        \zeta_{i,t}^{\mathrm{clip}}
        =
        \clip\!\left(\zeta_{i,t};-\zeta_{\max},\zeta_{\max}\right),
        \qquad
        \zeta_{\max}=20.
    \]
    The noise-trader quote is
    \[
        \tilde P_{i,j,t}^{N}=P_{j,t}\exp(\zeta_{i,t}^{\mathrm{clip}}).
    \]
    The order side follows the sign of the belief shock:
    \[
        \mathrm{side}_{i,j,t}^{N}=
        \begin{cases}
            \mathrm{Buy}, & \zeta_{i,t}>0,\\
            \mathrm{Sell}, & \zeta_{i,t}<0,\\
            \mathrm{Hold}, & \zeta_{i,t}=0.
        \end{cases}
    \]
    Noise-trader order size is bounded by a fraction of total mark-to-market
    wealth,
    \[
        B_{i,t}^{N}=\phi_N W_{i,t}.
    \]
    Given the quantised limit price \(P_{i,j,t}^{\text{lim}}\), submitted quantity is
    \[
        q_{i,j,t}^{N}=
        \begin{cases}
            \min\left\{\dfrac{B_{i,t}^{N}}{P_{i,j,t}^{\text{lim}}},
                    A_{i,t}^{B}(P_{i,j,t}^{\text{lim}})\right\}, & \mathrm{side}_{i,j,t}^{N}=\mathrm{Buy},\\[1.2ex]
            \min\left\{\dfrac{B_{i,t}^{N}}{P_{i,j,t}^{\text{lim}}},
            x_{i,j,t}+\dfrac{\bar w_N^{\text{short}}W_{i,t}}{P_{i,j,t}^{\text{lim}}}\right\}, & \mathrm{side}_{i,j,t}^{N}=\mathrm{Sell},\\[1.2ex]
            0, & \mathrm{side}_{i,j,t}^{N}=\mathrm{Hold}.
        \end{cases}
    \]
    Because only one asset is traded per step, the noise trader needs no running
    cash budget and uses the full balance \(A_{i,t}^{B}\).
    The second sell bound is a strategy-level short-exposure cap; the market
    coordinator still applies the global cash, inventory, margin, and
    borrowable-supply constraints before execution.

    \subsection{Buy-and-Hold Agents}

    Passive buy-and-hold agents \(i\in\mathcal{A}^H\) represent long-horizon
    investors who accumulate exposure toward a total target risky share
    \(\omega_H\). They are intentionally long-only and do not sell down after the
    target is reached. The target is split equally across assets rather than
    optimised, which keeps this population a passive liquidity benchmark rather
    than a second optimising agent class:
    \[
        x_{i,j,t}^{\ast,H}
        =
        \frac{W_{i,t}}{P_{j,t}}\cdot\frac{\omega_H}{n}.
    \]
    The submitted side is buy when current inventory is below this target and
    hold otherwise:
    \[
        \mathrm{side}_{i,j,t}^{H}
        =
        \begin{cases}
            \mathrm{Buy}, & x_{i,j,t}<x_{i,j,t}^{\ast,H},\\
            \mathrm{Hold}, & x_{i,j,t}\ge x_{i,j,t}^{\ast,H}.
        \end{cases}
    \]
    The quote is the latest market price, quantised to the buy side of that
    asset's price grid:
    \[
        P_{i,j,t}^{\text{lim},H}
        =
        \begin{cases}
            \max\{\delta_j,\delta_j\left\lfloor P_{j,t}/\delta_j\right\rfloor\}, & \delta_j>0,\\
            P_{j,t}, & \delta_j\le 0.
        \end{cases}
    \]
    The submitted quantity is therefore
    \[
        q_{i,j,t}^{H}
        =
        \begin{cases}
            \min\left\{\pospart{x_{i,j,t}^{\ast,H}-x_{i,j,t}},
            \tilde A_{i,t}^{B}(P_{i,j,t}^{\text{lim},H})\right\}, & \mathrm{side}_{i,j,t}^{H}=\mathrm{Buy},\\[1.2ex]
            0, & \mathrm{side}_{i,j,t}^{H}=\mathrm{Hold},
        \end{cases}
    \]
    again against the running cash budget \(\tilde c_{i,t}\).
    This population supplies persistent passive demand without adding a new
    behavioural forecast. It is useful for separating crisis dynamics generated by
    active behavioural rules from liquidity support supplied by long-horizon
    investors.

    \section{Valuation Quotes, Social Influence, and Limit Price Formation}\label{sec:valuation_quotes}

    The social module embeds agents in a synthetic two-dimensional information
    space. Each agent \(i\in\mathcal{A}\) is located at
    \((u_i,v_i)\in\mathbb{R}^2\) and, when social interaction is enabled, combines
    a private signal with lagged opinions transmitted by neighbouring agents. When
    social interaction is disabled, the implementation leaves opinions neutral, so
    the same valuation equations reduce to individual strategy quotes.

    \subsection{Spatial Structure and Neighbourhood}

    The Euclidean distance between agents \(i\) and \(l\) is defined as
    \begin{equation}
        d_{il} = \sqrt{{\left(u_i-u_l\right)}^2 + {\left(v_i-v_l\right)}^2}.
    \end{equation}

    Given an interaction radius \(r_{\mathrm{soc}}>0\), the neighbourhood of agent
    \(i\) is
    \begin{equation}
        \mathcal{N}_i = \{l \in \mathcal{A}\setminus\{i\} : d_{il} \le r_{\mathrm{soc}}\}.
    \end{equation}

    To capture local influence, we define a distance-based interaction weight
    \begin{equation}
        \omega_{il} = \frac{1}{1+d_{il}}, \qquad l \in \mathcal{N}_i.
    \end{equation}

    Hence, closer agents exert stronger influence on one another.

    \subsection{Private Signal}

    At each time step \(t\), agent \(i\) receives a private signal
    \begin{equation}
        \eta_{i,t} \sim \mathcal{N}(0,\sigma_{i,\text{priv}}^2),
    \end{equation}
    where \(\sigma_{i,\text{priv}}\) denotes the standard deviation of the private
    signal of agent \(i\).

    The private signal is interpreted as an idiosyncratic expected return signal or valuation shock.

    \subsection{Social Signal}

    Let \(o_{l,t-1}\) denote the opinion of agent \(l\) at time \(t-1\),
    interpreted as that agent's expected return signal. Agent \(i\) forms a social
    signal by aggregating the lagged opinions of neighbouring agents using
    distance-based weights:
    \begin{equation}
        S_{i,t}=
        \begin{cases}
            \displaystyle
            \frac{\sum_{l \in \mathcal{N}_i} \omega_{il} \, o_{l,t-1}}{\sum_{l \in \mathcal{N}_i} \omega_{il}}, & \text{if } |\mathcal{N}_i|>0,\\
            0, & \text{if } |\mathcal{N}_i|=0.
        \end{cases}
    \end{equation}

    \subsection{Opinion Formation}

    The opinion of agent \(i\) at time \(t\) is a convex combination of the
    private signal and the social signal:
    \begin{equation}
        o_{i,t} = \gamma_i S_{i,t} + (1-\gamma_i)\eta_{i,t},
        \qquad \gamma_i \in [0,1].
    \end{equation}

    The parameter \(\gamma_i\) measures the strength of social influence. If
    \(\gamma_i=0\), the agent relies entirely on private information. If
    \(\gamma_i=1\), the agent is fully socially driven.

    \subsection{Valuation Quote and Social Opinion Overlay}

    The active discretionary strategies submit valuation-based limit prices, one
    per asset they choose to trade. Social interaction, when enabled, shifts these
    quotes through the opinion \(o_{i,t}\). An agent holds a single opinion, which
    is applied identically to every asset it quotes: social influence is modelled
    as sentiment about the market, not as an asset-specific view. When social
    interaction is disabled, the code sets \(o_{i,t}=0\) for all agents, so the
    same equations reduce to individual valuation quotes without neighbour
    influence. Passive buy-and-hold agents use the market-anchored buy quote
    defined in the previous subsection. For chartists,
    \begin{equation}
        \tilde P_{i,j,t}^{C}=P_{j,t}\exp\!\left(\breve r_{i,j,t}^{C}+\kappa o_{i,t}\right),
    \end{equation}
    while for fundamentalists,
    \begin{equation}
        \tilde P_{i,j,t}^{F}=F_{j,t}\exp\!\left(\kappa o_{i,t}\right).
    \end{equation}
    Both valuation quotes are the price the agent expects at the end of its
    forecast horizon: the chartist's uses the whole horizon return
    \(\breve r_{i,j,t}^{C}\), and the fundamentalist's is the fundamental value
    the mispricing is expected to close to. The horizon therefore changes
    position sizes (Section~\ref{sec:allocation}) but not quotes.
    For noise traders, the stochastic quote
    \(\tilde P_{i,j,t}^{N}=P_{j,t}\exp(\zeta_{i,t}^{\mathrm{clip}})\) is generated
    directly from the random-liquidity belief shock rather than from the
    social-opinion overlay.
    Let \(\tilde P_{i,j,t}\) denote the applicable valuation quote, equal to
    \(\tilde P_{i,j,t}^{C}\) for chartists, \(\tilde P_{i,j,t}^{F}\) for
    fundamentalists, and \(\tilde P_{i,j,t}^{N}\) for noise traders.

    Define the effective log no-trade band
    \[
        \tau^{\mathrm{eff}}
        =
        \tau+\iota_{\mathrm{tc}}^D\tau_{\mathrm{tc}}^{\log}.
    \]
    For chartists and fundamentalists, no order is submitted in asset \(j\) when
    that asset's valuation quote remains sufficiently close to its market price:
    \begin{equation}
        \left|\log\frac{\tilde P_{i,j,t}}{P_{j,t}}\right|\le \tau^{\mathrm{eff}}.
    \end{equation}
    The test is applied asset by asset, so an agent may quote one asset and stay
    out of the other in the same step. Otherwise, the chartist or fundamentalist
    order side is
    \begin{equation}
        \mathrm{side}_{i,j,t}=
        \begin{cases}
            \mathrm{Buy}, & \tilde P_{i,j,t}>P_{j,t},\\
            \mathrm{Sell}, & \tilde P_{i,j,t}<P_{j,t}.
        \end{cases}
    \end{equation}
    For fundamentalists, forced liquidation in that asset overrides this rule and
    imposes a sell order. Noise traders instead use the sign rule in the preceding
    subsection and submit no order when the shock is exactly zero, the
    cost-adjusted log gap does not exceed \(\tau^{\mathrm{eff}}\), or feasible
    quantity is zero.

    If the asset's tick size is active, the submitted limit price is quantised on
    that asset's own price grid:
    \begin{equation}
        P_{i,j,t}^{\text{lim}}=
        \begin{cases}
            \max\{\delta_j,\delta_j\left\lfloor \tilde P_{i,j,t}/\delta_j\right\rfloor\}, & \mathrm{Buy},\\
            \max\{\delta_j,\delta_j\left\lceil \tilde P_{i,j,t}/\delta_j\right\rceil\}, & \mathrm{Sell}.
        \end{cases}
    \end{equation}
    For zero or numerically negligible tick sizes, the implementation leaves
    \(\tilde P_{i,j,t}\) unrounded. Otherwise, the same quantity rules as in the
    strategy definitions above are then used, replacing \(P_{j,t}\) by
    \(P_{i,j,t}^{\text{lim}}\) in the affordability constraint. After tick-size
    quantisation, discretionary chartist, fundamentalist, and noise-trader orders
    are again discarded when
    \[
        \left|\log\frac{P_{i,j,t}^{\text{lim}}}{P_{j,t}}\right|
        \le \tau^{\mathrm{eff}},
    \]
    because price rounding can move a marginal quote back inside the economically
    inactive band.

    \section{Continuous Limit Order Book and Price Update}\label{sec:continuous_lob}

    The current implementation uses a persistent continuous double-auction limit
    order book rather than a batch call auction. This moves price discovery from an
    auction-plus-impact equation to the interaction of bid-ask spread, visible
    depth, time priority, market orders, cancellations, and partial fills. The
    modelling choice is consistent with agent-based microstructure treatments of
    order-driven markets.

    \subsection{One Book per Asset}

    Each enabled asset owns one persistent order book, and books never cross
    assets. Every order carries an asset index \(\mathrm{asset}(\mathfrak{o})\)
    fixed at generation, which routes it to exactly one book; matching, lifecycle,
    depth, and price formation are then computed independently per asset. There is
    no cross-asset matching, no shared queue, and no cross-asset priority rule.
    Agents remain coupled across assets --- through shared cash, shared wealth,
    the joint Markowitz solve, and correlated fundamentals --- but the
    microstructure itself is separable. This separation is what makes each asset's
    price a genuine transaction price of its own market and lets stylised facts be
    screened per asset on identical terms.

    At the beginning of period \(t\), asset \(j\)'s book contains resting bids
    \(\mathcal{L}_{j,t}^{B}\) and resting asks \(\mathcal{L}_{j,t}^{A}\) carried
    over from previous periods. To avoid conflict with the social-opinion variable
    \(o_{i,t}\), an individual order record is denoted by \(\mathfrak{o}_{i,j,t}\)
    and represented as
    \[
        \mathfrak{o}_{i,j,t}=
        \left(
            j,\,
            i,\,
            \mathrm{side}_{i,j,t},\,
            \sigma_{i,t}^{\mathrm{ord}},\,
            P_{i,j,t}^{\text{lim}},\,
            q_{i,j,t},\,
            q_{i,j,t}^{\mathrm{rem}},\,
            T_{i,t}^{0},\,
            T_{i,t}^{X},\,
            n_{i,t}^{\mathrm{seq}}
        \right),
    \]
    where \(\sigma_{i,t}^{\mathrm{ord}}\in\{\mathrm{Limit},\mathrm{Market}\}\), the
    first component is the routing asset and the second the owning agent. The
    sequence number \(n_{i,t}^{\mathrm{seq}}\) is assigned by the book and is
    stable for the life of the order.

    The strategy layer still generates a side, quantity, and valuation price. The
    simulation coordinator converts an aggressive order to market style when
    \[
        \left|\log\frac{P_{i,j,t}^{\text{lim}}}{P_{j,t}}\right|
        \ge \theta_{\mathrm{MO}},
    \]
    or when a uniform random draw is below the configured probability
    \(\rho_{\mathrm{MO}}\). Limit-order residuals can rest in the book until they
    are executed, cancelled, expired, or invalidated by execution-time feasibility
    checks. Market-order residuals do not rest.

    \subsection{Order Lifecycle}

    Before new period-\(t\) submissions are collected, every book removes expired
    orders:
    \[
        \mathfrak{o}_{i,j,\tau}\ \text{expires at }t
        \quad\Longleftrightarrow\quad
        T_{i,\tau}^{X}>0\ \text{and}\ t>T_{i,\tau}^{X}.
    \]
    The book also applies endogenous stale-order cancellation. For a resting order
    with finite positive limit price, the cancellation probability is measured
    against its own asset's latest market price:
    \[
        \rho_{i,j,t}^{\mathrm{cancel}}
        =
        \clip\!\left(
                   \rho_0+\gamma_{\mathrm{mis}}
                   \left|\log\frac{P_{i,j,t}^{\text{lim}}}{P_{j,t}}\right|;
            0,1
        \right),
    \]
    where \(\rho_0\) is the base cancellation probability and
    \(\gamma_{\mathrm{mis}}\) controls sensitivity to distance from the latest
    market price. The draw is generated by the same seeded random engine as the
    rest of the simulation, so cancellations remain reproducible for a fixed
    configuration.

    Removals are counted in five separate categories per asset, because they carry
    different economic interpretations: stale cancellation (the rule above), order
    expiry, infeasible-resting cancellation and infeasible-incoming cancellation
    (Section~\ref{sec:short_sale_constraints}), self-trade prevention, and
    collar rejection. A legacy aggregate cancellation count is retained for
    continuity with earlier outputs.

    \subsection{Price-Time Priority}

    The rules below apply within one asset's book; \(j\) is fixed throughout this
    subsection. The bid book is sorted by descending price and then ascending
    sequence number:
    \[
        b\succ_B b'
        \quad\Longleftrightarrow\quad
        P_b^{\text{lim}}>P_{b'}^{\text{lim}}
        \ \text{or}\
        \left(P_b^{\text{lim}}=P_{b'}^{\text{lim}}
            \ \text{and}\ n_b^{\mathrm{seq}}<n_{b'}^{\mathrm{seq}}\right).
    \]
    The ask book is sorted by ascending price and then ascending sequence number:
    \[
        a\succ_A a'
        \quad\Longleftrightarrow\quad
        P_a^{\text{lim}}<P_{a'}^{\text{lim}}
        \ \text{or}\
        \left(P_a^{\text{lim}}=P_{a'}^{\text{lim}}
            \ \text{and}\ n_a^{\mathrm{seq}}<n_{a'}^{\mathrm{seq}}\right).
    \]
    Let \(a_1\) denote the best ask and \(b_1\) the best bid when the corresponding
    side is non-empty. A market buy crosses whenever \(a_1\) exists; a market sell
    crosses whenever \(b_1\) exists. Limit buys and sells cross only when
    \[
        P_{i,j,t}^{\text{lim}}\ge P_{a_1}^{\text{lim}},
        \qquad
        P_{i,j,t}^{\text{lim}}\le P_{b_1}^{\text{lim}},
    \]
    respectively. Non-crossing limit-order residuals are inserted into the
    appropriate side of the persistent book.

    \subsection{Self-Trade Prevention and the Market-Order Collar}\label{sec:stp_collar}

    Two execution-time screens are applied before a candidate match is priced.

    First, self-trade prevention. If the best opposing order belongs to the same
    agent as the incoming order,
    \[
        \mathrm{owner}(\mathfrak{o}^{\mathrm{rest}})
        =
        \mathrm{owner}(\mathfrak{o}^{\mathrm{in}}),
    \]
    no trade occurs. The resting order is removed from the book
    (cancel-resting semantics) and matching continues against the next-best
    counterparty. An agent trading with itself would otherwise create volume,
    transaction costs, and a printed transaction price out of a purely internal
    inventory shuffle, corrupting both the price series and the volume
    diagnostics. Cancelling the agent's own stale quote rather than the incoming
    order preserves the incoming order's ability to reach genuine counterparties.
    Removals of this kind are counted separately from stale cancellations.

    Second, the market-order price collar. For a market-style incoming order with
    candidate execution price \(p\), the trade is rejected when
    \[
        \left|\log\frac{p}{P_{j,t}}\right|>\theta_{\mathrm{col}},
        \qquad \theta_{\mathrm{col}}>0,
    \]
    where \(P_{j,t}\) is that asset's reference price at the start of the step. The
    incoming residual is then cancelled and the sweep stops. This is the standard
    exchange protection against a market order walking an empty book to an
    economically meaningless price. Setting \(\theta_{\mathrm{col}}=0\) disables
    the collar and restores unconstrained market-order sweeps, which is the
    baseline used in the pre-collar experiments.

    \subsection{Execution-Time Feasibility, Margin, and Borrow Constraints}\label{sec:short_sale_constraints}

    Cash, inventory, short-sale margin, and borrow constraints are revalidated at
    the moment of execution because a resting order may become infeasible after
    other trades change its owner's balance sheet. This is important in a
    persistent book: validity at order-submission time is not sufficient for
    validity at later execution time.

    Borrow constraints are per asset: shares of asset \(j\) can only be borrowed
    against asset \(j\)'s own lending base. Let
    \[
        \bar X_{j,0}=N\max\{x_{j,0},0\}
    \]
    denote asset \(j\)'s initial share supply used as its securities-lending base.
    At a candidate execution price \(p>0\) in asset \(j\), the equity available to
    the margin rule is
    \[
        E_{i,j,t}(p)=c_{i,t}+x_{i,j,t}p.
    \]
    This is the agent's shared cash plus the position in the asset being traded;
    positions in other assets do not enter. Margin is therefore assessed asset by
    asset with no cross-asset netting, which is a deliberate simplification
    relative to a portfolio-margin rule and is conservative in the sense that it
    never lets a gain in one asset finance short exposure in another.

    The initial-margin rule permits a maximum short position
    \[
        s_{i,j,t}^{\max}(p)
        =
        \max\left\{0,\frac{E_{i,j,t}(p)}{M_0^{S}p}\right\}.
    \]
    The agent-level capacity to increase short exposure in that asset is therefore
    \[
        A_{i,j,t}^{S}(p)=\pospart{s_{i,j,t}^{\max}(p)-\pospart{-x_{i,j,t}}}.
    \]
    The borrow pool of asset \(j\) remaining at the instant of validation is
    \[
        B_{j,t}^{\mathrm{rem}}
        =
        \pospart{\varphi^{B}\bar X_{j,0}-\mathrm{SI}_{j,t}},
        \qquad
        \mathrm{SI}_{j,t}=\sum_{l\in\mathcal{A}}\pospart{-x_{l,j,t}}.
    \]
    For buyer \(b\), seller \(s\), resting order \(r\), incoming order \(u\), all in
    asset \(j\), and execution price \(p=P_{r}^{\text{lim}}\), the executable
    matched quantity is
    \[
        m
        =
        \min\left\{
                q_{u}^{\mathrm{rem}},
            q_{r}^{\mathrm{rem}},
                A_{b,t}^{B}(p),
                \bar q_{s,j,t}^{S}(p)
        \right\},
    \]
    where the buyer-side bound \(A_{b,t}^{B}(p)=c_{b,t}/\bigl(p(1+f_{\mathrm{tc}})\bigr)\)
    uses the buyer's \emph{actual} cash balance --- the running submission-time
    budget of Section~\ref{sec:allocation} constrains order generation only --- and
    the seller-side feasible quantity is
    \[
        \bar q_{s,j,t}^{S}(p)=
        \begin{cases}
            \pospart{x_{s,j,t}}, & \text{short sales disabled},\\[1ex]
            \pospart{x_{s,j,t}}
            +
            \min\left\{
                    \pospart{q_{s}^{\mathrm{rem}}-\pospart{x_{s,j,t}}},\,
                    A_{s,j,t}^{S}(p),\,
                    B_{j,t}^{\mathrm{rem}}
            \right\}, & \text{short sales enabled}.
        \end{cases}
    \]
    A candidate match must also clear a minimum economic fill before it is
    executed. Let the requested match quantity be
    \[
        m^{\mathrm{req}}
        =
        \min\left\{q_{u}^{\mathrm{rem}},\,q_{r}^{\mathrm{rem}}\right\},
    \]
    and define the minimum economic fill
    \[
        m^{\min}
        =
        \max\left\{\epsilon_q,\ \varphi_{\mathrm{fill}}\,m^{\mathrm{req}}\right\},
        \qquad
        \epsilon_q=10^{-8},
        \qquad
        \varphi_{\mathrm{fill}}=10^{-6}.
    \]
    The trade is executed only when \(m>m^{\min}\). Otherwise the fill is treated
    as economically negligible and the binding side is cancelled: the buy side
    when \(A_{b,t}^{B}(p)\le m^{\min}\) and the sell side when
    \(\bar q_{s,j,t}^{S}(p)\le m^{\min}\). Cancelling removes the resting order
    from the book when the resting side is infeasible and discards the incoming
    residual when the incoming side is infeasible. The relative component
    \(\varphi_{\mathrm{fill}}m^{\mathrm{req}}\) addresses a degeneracy in which
    a binding constraint permits a strictly positive but negligible fill. In
    particular, when an asset's short interest sits exactly at its
    borrowable-supply cap and a covering buyer meets a shorting seller, each
    executed share releases exactly as much borrow capacity as it consumes, so
    the permitted fill never grows and the matching loop would otherwise grind
    a single order through economically meaningless micro-fills.

    As a terminal numerical safeguard, the matching loop for one incoming order
    is bounded by \(K_{\mathrm{match}}=10^{5}\) iterations. If the bound is
    reached, the incoming order's residual is cancelled and a warning is logged.
    Legitimate matching touches at most a few price levels per order, so this
    bound is inactive except under unforeseen constraint interactions.

    These
    checks operationalise institutional short-sale frictions: pessimistic agents
    can transmit negative beliefs only when they have sufficient equity and when
    borrowable shares remain available.

    \subsection{Matching and Accounting}

    Orders are processed sequentially in the randomly shuffled activation order.
    An incoming order repeatedly trades against the best opposite order while the
    crossing condition holds, both sides remain feasible, and the executable
    quantity exceeds the minimum economic fill \(m^{\min}\). The execution price is
    the resting order's limit price:
    \[
        p_k=P_{r_k}^{\text{lim}}.
    \]
    For each execution of \(m_k>m_k^{\min}\) shares between buyer \(b\) and seller \(s\),
    define the gross notional and one-side transaction cost
    \[
        G_k=m_kp_k,
        \qquad
        F_k^{\mathrm{tc}}=f_{\mathrm{tc}}G_k.
    \]
    Portfolios update immediately. Cash is shared, while inventory moves only in
    the traded asset \(j\):
    \[
        c_{b,t}^{+}=c_{b,t}-G_k-F_k^{\mathrm{tc}},
        \qquad
        x_{b,j,t}^{+}=x_{b,j,t}+m_k,
    \]
    \[
        c_{s,t}^{+}=c_{s,t}+G_k-F_k^{\mathrm{tc}},
        \qquad
        x_{s,j,t}^{+}=x_{s,j,t}-m_k.
    \]
    The implementation also updates each filled agent's trade count, traded
    volume, per-asset current-position average entry price, per-asset short-sale
    state, cumulative transaction costs,
    \[
        \Gamma_{b,t}^{+}=\Gamma_{b,t}+F_k^{\mathrm{tc}},
        \qquad
        \Gamma_{s,t}^{+}=\Gamma_{s,t}+F_k^{\mathrm{tc}},
    \]
    and trade-ledger records at the same instant. Partial fills reduce the remaining quantity of
    both orders:
    \[
        q_{u}^{\mathrm{rem}}\leftarrow q_{u}^{\mathrm{rem}}-m_k,
        \qquad
        q_{r}^{\mathrm{rem}}\leftarrow q_{r}^{\mathrm{rem}}-m_k.
    \]
    The resting order is removed when \(q_{r}^{\mathrm{rem}}\) reaches zero. A
    limit incoming order with positive non-crossing residual rests in the book
    after being clipped to its owner's current feasibility constraint. A market
    incoming order never rests.

    The endogenous executed volume, endogenous transaction-cost total, and last
    transaction price of asset \(j\) in period \(t\) are
    \[
        Q_{j,t}^{\mathrm{exec}}=\sum_{k\in\mathcal{K}_{j,t}} m_k,
        \qquad
        \Theta_{j,t}^{\mathrm{exec}}=2\sum_{k\in\mathcal{K}_{j,t}} F_k^{\mathrm{tc}},
        \qquad
        P_{j,t}^{\mathrm{last}}=
        \begin{cases}
            p_{K_j}, & Q_{j,t}^{\mathrm{exec}}>0,\\
            P_{j,t}, & Q_{j,t}^{\mathrm{exec}}=0,
        \end{cases}
    \]
    where \(\mathcal{K}_{j,t}\) is the set of asset-\(j\) executions in the period
    and \(p_{K_j}\) is the final one observed.

    The same endogenous continuous-double-auction executions define a gross
    Smith-style allocative-efficiency diagnostic. For this diagnostic, the limit
    price attached to each submitted or resting order is treated as the observable
    reservation-value proxy. This is consistent with the current strategy layer,
    where quotes are generated from valuation-based beliefs before an order is
    optionally converted to market style. Let \(v_{o,t}\) be this order price and
    let \(\mathcal{K}_{j,t}^{\mathrm{LOB}}\) be the set of endogenous LOB
    executions in asset \(j\) in period \(t\). Forced buy-ins are excluded because
    they are broker risk-control trades rather than voluntary double-auction
    allocations. The diagnostic is computed independently for each asset, since
    surplus from one book is not comparable with surplus from another. The
    realised gross allocative surplus is
    \[
        AS_{j,t}
        =
        \sum_{k\in\mathcal{K}_{j,t}^{\mathrm{LOB}}}
        m_k\left(v_{b_k,t}-v_{s_k,t}\right),
    \]
    where \(b_k\) and \(s_k\) are the buyer and seller orders matched in
    execution \(k\). Transaction costs are deliberately not subtracted here; they
    remain in \(\Theta_{j,t}^{\mathrm{exec}}\) as a separate institutional-friction
    diagnostic.

    The maximum feasible gross surplus is computed from asset \(j\)'s surviving
    resting book after lifecycle expiry and cancellation plus all new submitted
    orders for that asset before matching. Denote these buy and sell order
    multisets by \(\mathcal{O}_{j,t}^B\) and \(\mathcal{O}_{j,t}^S\). Sort buys
    from highest to lowest reservation value and sells from lowest to highest
    reservation value. Greedily match marginal quantities while the buyer
    reservation value exceeds the seller reservation value:
    \[
        MS_{j,t}
        =
        \sum_{\iota}
        m_{j,\iota,t}^{\ast}
        \left(v_{j,\iota,t}^{B,\ast}-v_{j,\iota,t}^{S,\ast}\right),
        \qquad
        v_{j,\iota,t}^{B,\ast}>v_{j,\iota,t}^{S,\ast}.
    \]
    The per-step allocative-efficiency ratio is therefore
    \[
        AE_{j,t}=
        \begin{cases}
            AS_{j,t}/MS_{j,t}, & MS_{j,t}>0,\\
            0, & MS_{j,t}=0.
        \end{cases}
    \]
    Steps with \(MS_{j,t}=0\) have no positive-surplus double-auction opportunity
    and are excluded from that asset's run-level mean:
    \[
        \overline{AE}_j
        =
        \frac{1}{|\mathcal{T}_{AE,j}|}
        \sum_{t\in\mathcal{T}_{AE,j}} AE_{j,t},
        \qquad
        \mathcal{T}_{AE,j}=\{t:MS_{j,t}>0\}.
    \]

    \subsection{Spread, Depth, and Price Formation}

    After matching, book diagnostics are computed per asset from the remaining
    books:
    \[
        P_{j,t}^B=
        \begin{cases}
            \max_{b\in\mathcal{L}_{j,t}^B} P_b^{\text{lim}}, & \mathcal{L}_{j,t}^B\ne\varnothing,\\
            0, & \mathcal{L}_{j,t}^B=\varnothing,
        \end{cases}
        \qquad
        P_{j,t}^A=
        \begin{cases}
            \min_{a\in\mathcal{L}_{j,t}^A} P_a^{\text{lim}}, & \mathcal{L}_{j,t}^A\ne\varnothing,\\
            0, & \mathcal{L}_{j,t}^A=\varnothing.
        \end{cases}
    \]
    When both quotes exist,
    \[
        S_{j,t}^{BA}=P_{j,t}^A-P_{j,t}^B,
        \qquad
        M_{j,t}^Q=\frac{P_{j,t}^A+P_{j,t}^B}{2}.
    \]
    Visible depth \(D_{j,t}^B(h)\) and \(D_{j,t}^A(h)\) are the sums of remaining
    quantity across the first \(h\) distinct price levels on the bid and ask sides.

    The model records submitted net demand and notional imbalance before execution.
    These aggregates are reported for the primary asset in the legacy single-asset
    view:
    \[
        \mathrm{ND}_t=
        \sum_{i:\mathrm{Buy}} q_{i,j^{\ast},t}
        -
        \sum_{i:\mathrm{Sell}} q_{i,j^{\ast},t},
    \]
    \[
        \mathrm{NI}_t=\mathrm{ND}_t\,P_{j^{\ast},t},
        \qquad
        \bar I_t=\clip\!\left(\frac{\mathrm{NI}_t}{C_0};-1,1\right),
    \]
    where \(C_0=N\max\{c_0,1\}\). It also records the relative fundamental signal
    \[
        g_t^{F}=\frac{F_{j^{\ast},t+1}-P_{j^{\ast},t}}{\max\{P_{j^{\ast},t},\epsilon_P\}}.
    \]
    These variables remain useful diagnostics for order-flow pressure and
    valuation pressure, but in continuous LOB mode they do not directly move the
    stored market price. Their economic effect enters through future agent orders.
    The legacy parameters \(\beta_F\), \(\beta_D\), and \(\sigma_P\) are therefore
    inactive under this microstructure.

    Price formation is resolved independently for each asset. The trade-execution
    indicator is
    \[
        \chi_{j,t}=\mathbf{1}\!\left(Q_{j,t}^{\mathrm{exec}}>0\right),
    \]
    and the official market price is the last transaction price when a trade
    occurs in that asset and the previous market price otherwise:
    \[
        P_{j,t+1}=
        \begin{cases}
            \max\{\epsilon_P,P_{j,t}^{\mathrm{last}}\}, & \chi_{j,t}=1,\\
            P_{j,t}, & \chi_{j,t}=0.
        \end{cases}
    \]
    The realised simple and log returns are
    \[
        r_{j,t}^{\mathrm{raw}}=
        \begin{cases}
            \dfrac{P_{j,t+1}-P_{j,t}}{\max\{P_{j,t},\epsilon_P\}}, & \chi_{j,t}=1,\\[1.4ex]
            0, & \chi_{j,t}=0,
        \end{cases}
        \qquad
        \ell_{j,t+1}=\log\frac{P_{j,t+1}}{P_{j,t}}.
    \]
    Thus a no-trade step always carries that asset's price forward and has zero
    realised log return, and one asset can trade in a step while the other does
    not. Price changes are transaction-price changes, not exogenous impact
    updates.

    \subsection{Cross-Sectional Portfolio-Weight Diagnostic}

    To make cross-asset allocation observable at run time, each step records the
    cross-sectional mean weight held in each asset. For agent \(i\) with strictly
    positive wealth,
    \[
        w_{i,j,t}^{\mathrm{held}}
        =
        \frac{x_{i,j,t}P_{j,t}}{W_{i,t}},
        \qquad
        \bar w_{j,t}
        =
        \frac{1}{|\mathcal{A}_t^{+}|}
        \sum_{i\in\mathcal{A}_t^{+}} w_{i,j,t}^{\mathrm{held}},
        \qquad
        \mathcal{A}_t^{+}=\{i:W_{i,t}>0\}.
    \]
    Agents with non-positive wealth are excluded because their weights are not
    economically interpretable and would dominate the mean. Weights are not
    clipped: leverage and short positions legitimately push individual weights
    outside \([0,1]\). Because the mean is linear, the mean cash weight is
    recovered downstream as \(1-\sum_j\bar w_{j,t}\) and needs no column of its
    own.

    \subsection{Maintenance Margin and Forced Buy-in}

    After every asset's market price \(P_{j,t+1}\) has been stored, the
    implementation marks all short positions to market and applies the
    maintenance-margin rule. The rule is evaluated asset by asset, in declared
    index order, against the agent's shared cash. There is no cross-asset netting,
    so sequential enforcement may over-cover relative to a portfolio-margin rule;
    this is a documented modelling simplification consistent with the
    execution-time margin treatment in Section~\ref{sec:short_sale_constraints}.

    Let \((c_{i,t}^{L},x_{i,j,t}^{L})\) denote the portfolio after the endogenous
    LOB execution of period \(t\) and before any forced buy-in, and let
    \(c_{i,j,t}^{L}\) denote the cash balance actually available when asset \(j\)
    is processed: it equals \(c_{i,t}^{L}\) for the first asset in index order and
    is reduced by each earlier asset's cover. Define
    \[
        E_{i,j,t}^{L}=c_{i,j,t}^{L}+x_{i,j,t}^{L}P_{j,t+1},
        \qquad
        s_{i,j,t}^{L}=\pospart{-x_{i,j,t}^{L}}.
    \]
    The post-LOB short margin ratio in asset \(j\) is
    \[
        \mathcal{M}_{i,j,t}^{L}=
        \begin{cases}
            \infty, & s_{i,j,t}^{L}=0,\\[1ex]
            \dfrac{E_{i,j,t}^{L}}{s_{i,j,t}^{L}P_{j,t+1}}, & s_{i,j,t}^{L}>0.
        \end{cases}
    \]
    If \(\mathcal{M}_{i,j,t}^{L}\ge M_{\min}^{S}\), no forced buy-in occurs. If
    \(\mathcal{M}_{i,j,t}^{L}<M_{\min}^{S}\), the broker-style risk control reduces
    the short position enough to restore the stricter initial-margin ratio when
    cash permits:
    \[
        b_{i,j,t}^{\mathrm{req}}
        =
        \pospart{
            s_{i,j,t}^{L}
            -
            \frac{E_{i,j,t}^{L}}{M_0^{S}P_{j,t+1}}
        }.
    \]
    Let the effective forced-buy-in transaction-cost rate be
    \[
        f_{\mathrm{tc}}^{B}=\iota_{\mathrm{tc}}^B f_{\mathrm{tc}}.
    \]
    The actual forced buy-in quantity is clipped by the current short position and
    the cash still available after any earlier asset's cover in the same step,
    net of any enabled forced-buy-in cost:
    \[
        b_{i,j,t}^{\text{cover}}
        =
        \begin{cases}
            0, & \mathcal{M}_{i,j,t}^{L}\ge M_{\min}^{S},\\[1ex]
            \min\left\{
                    b_{i,j,t}^{\mathrm{req}},\,
                    s_{i,j,t}^{L},\,
                    \dfrac{c_{i,j,t}^{L}}{P_{j,t+1}(1+f_{\mathrm{tc}}^{B})}
            \right\}, & \mathcal{M}_{i,j,t}^{L}<M_{\min}^{S}.
        \end{cases}
    \]
    The forced cover is executed against an external liquidity account at the
    stored market price \(P_{j,t+1}\). Its notional and fee are
    \[
        G_{i,j,t}^{B}=b_{i,j,t}^{\text{cover}}P_{j,t+1},
        \qquad
        F_{i,j,t}^{B}=f_{\mathrm{tc}}^{B}G_{i,j,t}^{B}.
    \]
    The post-buy-in portfolio is
    \[
        c_{i,t+1}=c_{i,t}^{L}-\sum_{j}\left(G_{i,j,t}^{B}+F_{i,j,t}^{B}\right),
        \qquad
        x_{i,j,t+1}=x_{i,j,t}^{L}+b_{i,j,t}^{\text{cover}}.
    \]
    Cumulative transaction costs are increased by \(\sum_j F_{i,j,t}^{B}\) for the
    covered agent.
    Forced buy-ins are recorded in the trade ledger with an external counterparty
    identifier and are attributed to the asset they cover in reported traded
    volume:
    \[
        \widetilde Q_{j,t}
        =
        Q_{j,t}^{\text{exec}}
        +
        \sum_{i\in\mathcal{A}} b_{i,j,t}^{\text{cover}}.
    \]
    The period transaction-cost diagnostic for asset \(j\) is
    \[
        \Theta_{j,t}=\Theta_{j,t}^{\mathrm{exec}}+\sum_{i\in\mathcal{A}}F_{i,j,t}^{B}.
    \]
    They do not alter the price-formation indicator \(\chi_{j,t}\), because they
    are applied after the market price for the period has already been formed. This
    choice separates behavioural order-flow price discovery from broker-enforced
    risk reduction while still allowing short squeezes to appear in agent balance
    sheets and volume diagnostics.

    \section{Cumulative Prospect Theory Preference Overlay}\label{sec:cpt}

    CPT is implemented as a behavioural preference transform applied to the
    strategy-specific expected excess return. It does not create a separate trading
    signal. Instead, it converts the baseline forecast \(\hat r_{i,j,t}\) into the
    allocation input \(\breve r_{i,j,t}\), the certainty-equivalent premium of
    holding the asset for one step over holding cash
    (Section~\ref{sec:cpt_premium}), for the seeded fraction
    \(\varphi_{\mathrm{CPT}}\) of agents assigned to CPT\@. Agents not assigned to
    CPT use \(\breve r_{i,j,t}=\hat r_{i,j,t}\). This follows the
    behavioural-finance interpretation of Prospect Theory and rank-dependent
    Cumulative Prospect Theory: agents evaluate outcomes as gains or losses
    relative to a reference point, penalise losses more heavily than gains, and
    distort cumulative probabilities. The transform is applied separately to each
    asset's forecast, and the resulting vector enters the joint Markowitz solve of
    Section~\ref{sec:allocation}.

    The entire valuation is carried out in \emph{return space}. Both the scenario
    outcomes and the reference point are log returns, so the transform is
    invariant to the price and wealth level of the underlying asset: two assets
    with identical forecasts and identical volatilities receive identical
    prospect-adjusted returns even when one trades at \(100\) and the other at
    \(40\). Under the earlier price-space formulation the same forecast produced
    different behavioural adjustments at different price levels, which would have
    made a multi-asset CPT population incoherent.

    \subsection{Value Function}

    Let \(x\) be a gain or loss in return units, measured relative to the
    reference log return \(\varrho_{i,j,t}\). Agent \(i\)'s value function is
    \[
        v_i(x)=
        \begin{cases}
            x^{\alpha_i^G}, & x\ge 0,\\[0.8ex]
            -\lambda_i{\left(-x\right)}^{\alpha_i^L}, & x<0,
        \end{cases}
    \]
    where \(\lambda_i>0\) controls loss aversion and
    \(\alpha_i^G,\alpha_i^L>0\) control diminishing sensitivity in the gain and
    loss domains. The canonical empirical calibration is
    \(\lambda=2.25\) and \(\alpha^G=\alpha^L=0.88\). The working configuration
    retains the curvature values \(\alpha^G=\alpha^L=0.88\) but treats the
    population loss-aversion level and its cross-sectional dispersion as
    calibrated parameters, so the configured \(\lambda\) is set by the
    calibration harness rather than fixed at the canonical value. Because the
    transform reports a certainty equivalent in log-return units and its
    scenario ladder is moment matched (Sections~\ref{sec:cpt_scenarios}
    and~\ref{sec:cpt_premium}), a calibrated \(\lambda\) is directly comparable
    with the empirical values; values calibrated before that correction are not.

    \subsection{Probability Weighting}

    For domain \(d\in\{G,L\}\), the decision-weighting function is
    \[
        w_i^d(p)=
        \frac{p^{\psi_i^d}}
        {{\left(p^{\psi_i^d}+\bar p^{\psi_i^d}\right)}^{1/\psi_i^d}},
        \qquad
        \bar p=1-p,
        \qquad p\in[0,1],
    \]
    with \(w_i^d(0)=0\) and \(w_i^d(1)=1\). The gain and loss probability
    distortion parameters are \(\psi_i^G\) and \(\psi_i^L\), respectively.

    \subsection{Rank-Dependent Decision Weights}

    Consider a finite scenario set for asset \(j\),
    \[
        \mathcal{X}_{i,j,t}=\{(x_{i,j,t,s},p_s):s=1,\ldots,n_{\mathrm{CPT}}\},
        \qquad
        \sum_s p_s=1.
    \]
    Partition the scenarios into gains \(x_{i,j,t,s}\ge 0\) and losses
    \(x_{i,j,t,s}<0\). Gains are sorted from best to worst, while losses are sorted
    from worst to best. For the sorted gain outcomes, define
    \[
        \pi_{i,k}^{G}
        =
        w_i^G\!\left(\sum_{h=1}^{k}p_h^{G}\right)
        -
        w_i^G\!\left(\sum_{h=1}^{k-1}p_h^{G}\right).
    \]
    For the sorted loss outcomes, define
    \[
        \pi_{i,k}^{L}
        =
        w_i^L\!\left(\sum_{h=1}^{k}p_h^{L}\right)
        -
        w_i^L\!\left(\sum_{h=1}^{k-1}p_h^{L}\right).
    \]
    The CPT value of the scenario set is
    \[
        \mathcal{V}_{i,j,t}
        =
        \sum_{k\in G}\pi_{i,k}^{G}v_i(x_{i,j,t,k}^{G})
        +
        \sum_{k\in L}\pi_{i,k}^{L}v_i(x_{i,j,t,k}^{L}).
    \]
    This rank-dependent construction is important for crisis experiments because it
    allows adverse tail scenarios to receive disproportionate decision weight.

    \subsection{Return Scenario Construction}\label{sec:cpt_scenarios}

    CPT values the trade over the agent's forecast horizon \(\mathcal H_i\)
    (Section~\ref{sec:allocation}), which is therefore also its evaluation period
    in the sense of Benartzi and Thaler. Over that horizon the forecast
    \(\hat r_{i,j,t}\) is the expected excess log return, cash earns
    \(\mathcal H_i\ell_f\), and, with independent per-step returns, the volatility
    is
    \[
        \sigma_{i,j,t}^{\mathcal H}=\sigma_{j,t}\sqrt{\mathcal H_i}.
    \]
    Judging the trade over a single step instead would set a small per-step
    forecast against a full step of volatility, and loss aversion would then rule
    out holding the asset at almost any plausible mispricing; this is Benartzi and
    Thaler's myopic loss aversion. With \(\mathcal H_i=1\) everything below
    reduces to the one-step valuation.

    Let \(n_{\mathrm{CPT}}=2H+1\) be an odd integer greater than or equal to
    three. The ladder represents the agent's subjective distribution of the log
    return over the horizon, centred on \(\mathcal H_i\ell_f+\hat r_{i,j,t}\) with
    standard deviation \(\sigma_{i,j,t}^{\mathcal H}\). For scenario index
    \(s=0,\ldots,2H\), define the binomial probability and the moment-matched step
    \[
        p_s=\binom{2H}{s}2^{-2H},
        \qquad
        \Delta_{i,j,t}=\frac{2\,\clip(\sigma_{i,j,t}^{\mathcal H};0,0.95)}{\sqrt{n_{\mathrm{CPT}}-1}},
    \]
    and the subjective return scenario
    \[
        R_{i,j,t,s}^{\mathrm{sc}}
        =
        \clip\!\left(
                   \mathcal H_i\ell_f+\clip(\hat r_{i,j,t};-0.95,0.95)
                   +(s-H)\,\Delta_{i,j,t};
            -0.95,0.95
        \right).
    \]
    Because \(\operatorname{Var}(s-H)=H/2\), in the absence of clipping the
    ladder has mean \(\mathcal H_i\ell_f+\hat r_{i,j,t}\) and variance
    \((\sigma_{i,j,t}^{\mathcal H})^2\) for every \(n_{\mathrm{CPT}}\), and it
    converges in distribution to the normal distribution with these moments as
    \(n_{\mathrm{CPT}}\) grows (de Moivre--Laplace). This is the variance-matching
    rule of the Cox--Ross--Rubinstein binomial tree. The scenario count therefore
    controls only the resolution of the tails, never the risk the agent
    perceives. An unscaled step \(\Delta=\sigma\) would give the perceived
    variance \(\sigma^2(n_{\mathrm{CPT}}-1)/4\), letting a numerical setting act as
    a preference parameter. When \(n_{\mathrm{CPT}}=3\), the scenario
    probabilities are \(\{0.25,0.50,0.25\}\), placed at
    \(\mathcal H_i\ell_f+\hat r_{i,j,t}\pm\sqrt2\,\sigma_{i,j,t}^{\mathcal H}\) and
    \(\mathcal H_i\ell_f+\hat r_{i,j,t}\).

    The corresponding gain or loss is the scenario return measured relative to the
    reference log return, so it is itself a return:
    \[
        x_{i,j,t,s}=R_{i,j,t,s}^{\mathrm{sc}}-\varrho_{i,j,t}.
    \]

    \subsection{Certainty Equivalent and Comparison with Cash}\label{sec:cpt_premium}

    CPT is a theory of choice between prospects, so the value of holding the asset
    is meaningful only against the agent's alternative. Over the horizon the
    alternative is to sell and hold cash, whose sure outcome relative to the same
    reference realises today's paper gain or loss:
    \[
        x_{i,j,t}^{\mathrm{c}}=\mathcal H_i\ell_f-\varrho_{i,j,t}.
    \]
    Applying the value functional of the previous subsections to the scenario
    ladder and to the sure outcome gives
    \[
        \mathcal{V}_{i,j,t}^{\mathrm{hold}}
        =
        \mathcal{V}\bigl(\{(x_{i,j,t,s},p_s)\}_{s=0}^{2H}\bigr),
        \qquad
        \mathcal{V}_{i,j,t}^{\mathrm{cash}}=v_i\bigl(x_{i,j,t}^{\mathrm{c}}\bigr),
    \]
    the latter because a sure outcome carries the single decision weight
    \(w(1)-w(0)=1\). Prospect values are measured in value units rather than
    returns, so both are mapped back through the inverse value function
    \[
        v_i^{-1}(y)=
        \begin{cases}
            y^{1/\alpha_i^G}, & y\ge 0,\\[0.8ex]
            -\left(-y/\lambda_i\right)^{1/\alpha_i^L}, & y<0,
        \end{cases}
    \]
    which yields each option's certainty equivalent: the sure log return over the
    horizon, measured from today's price, that the agent values as highly as the
    option,
    \[
        \mathrm{CE}_{i,j,t}^{\mathrm{hold}}
        =\varrho_{i,j,t}+v_i^{-1}\bigl(\mathcal{V}_{i,j,t}^{\mathrm{hold}}\bigr),
        \qquad
        \mathrm{CE}_{i,j,t}^{\mathrm{cash}}
        =\varrho_{i,j,t}+v_i^{-1}\bigl(\mathcal{V}_{i,j,t}^{\mathrm{cash}}\bigr)=\mathcal H_i\ell_f.
    \]
    Tversky and Kahneman measured CPT preferences through such certainty
    equivalents. The prospect-adjusted excess return is the certainty-equivalent
    premium of holding over cash,
    \[
        \mathcal{T}_i(\hat r_{i,j,t},\sigma_{j,t};\varrho_{i,j,t})
        =
        \mathrm{CE}_{i,j,t}^{\mathrm{hold}}-\mathrm{CE}_{i,j,t}^{\mathrm{cash}}
        =
        v_i^{-1}\bigl(\mathcal{V}_{i,j,t}^{\mathrm{hold}}\bigr)+\varrho_{i,j,t}-\mathcal H_i\ell_f,
    \]
    \[
        \breve r_{i,j,t}
        =
        \clip\!\left(
                   \mathcal{T}_i(\hat r_{i,j,t},\sigma_{j,t};\varrho_{i,j,t});
            -1,1
        \right),
    \]
    a return over the horizon, which the allocation rule of
    Section~\ref{sec:allocation} divides by \(\mathcal H_i\).

    The construction has four properties, each pinned by a unit test. First,
    \(v_i\) is strictly increasing, so \(\breve r_{i,j,t}>0\) exactly when
    \(\mathcal{V}_{i,j,t}^{\mathrm{hold}}>\mathcal{V}_{i,j,t}^{\mathrm{cash}}\):
    the sign of the premium is the hold-versus-sell choice. Second, a riskless
    forecast passes through unchanged: \(\sigma_{j,t}=0\) gives
    \(\breve r_{i,j,t}=\hat r_{i,j,t}\) for any reference point, so CPT alters how
    agents treat risk rather than how they read a sure return. Third, with
    \(\alpha_i^G=\alpha_i^L\) the premium is homogeneous of degree one in
    \((\hat r,\sigma^{\mathcal H},\varrho,\mathcal H_i\ell_f)\), so it is measured
    in return units commensurate with the covariance in the Markowitz rule and
    does not depend on the time unit of a step. Fourth, if a paper loss or gain of
    size \(\delta\) satisfies
    \[
        \delta>|\mathcal H_i\ell_f+\hat r_{i,j,t}|+\sigma_{i,j,t}^{\mathcal H}\sqrt{n_{\mathrm{CPT}}-1},
    \]
    every scenario is a loss when \(\varrho_{i,j,t}=+\delta\) and a gain when
    \(\varrho_{i,j,t}=-\delta\). The convex loss branch then makes the agent risk
    seeking and the concave gain branch risk averse, the reflection effect of
    Kahneman and Tversky, so
    \[
        \mathcal{T}_i(\hat r_{i,j,t},\sigma_{j,t};+\delta)
        >\hat r_{i,j,t}>
        \mathcal{T}_i(\hat r_{i,j,t},\sigma_{j,t};-\delta).
    \]
    An agent holding a loser is therefore more willing to keep it than an agent
    holding a winner, the direction of the disposition effect. When the gamble
    straddles the reference the ranking is not guaranteed, consistent with the
    finding that prospect theory predicts the disposition effect only
    conditionally.

    The premium enters the mean-variance rule of Section~\ref{sec:allocation} as
    an expected excess return over the horizon, so risk is priced twice: by CPT
    through the premium and by the allocation rule through \(\rho_i\Sigma_t\).
    This hybrid is a deliberate simplification. Optimising the CPT value directly
    over portfolio weights is not used because, with power value functions, that
    problem is ill-posed or yields corner solutions.

    \subsection{Agent-Specific Reference Points}\label{sec:cpt_reference_points}

    The implementation resolves the CPT anchor in two stages: it first selects a
    behavioural anchor \(\nu_{i,j,t}\) in price or wealth units, then maps that
    anchor into the reference log return \(\varrho_{i,j,t}\) actually used by
    \(\mathcal{T}_i\). The configured mode
    \[
        m_{\nu}\in\{
        \mathrm{configured},
        \mathrm{entry},
        \mathrm{peakwealth},
        \mathrm{movingaverage},
        \mathrm{initialwealth}
        \}
    \]
    selects the behavioural anchor. These compact mathematical labels correspond
    to the configuration values \texttt{configured},
    \texttt{average\_purchase\_price}, \texttt{recent\_peak\_wealth},
    \texttt{moving\_average\_price}, and \texttt{initial\_wealth}. The fixed
    configured reference point \(\nu_i^0>0\) is retained as the fallback and
    reproduces the original fixed-reference specification when
    \(m_{\nu}=\mathrm{configured}\).

    The personal position reference is updated from executed trades, not submitted
    orders, and is maintained separately for each asset. Let \(e_{i,j,t}^{+}\)
    denote an executed buy quantity and \(e_{i,j,t}^{-}\) an executed sell quantity
    in asset \(j\). The signed execution quantity is
    \[
        \Delta x_{i,j,t}^{\mathrm{exec}}=e_{i,j,t}^{+}-e_{i,j,t}^{-}.
    \]
    Before the inventory update is stored, the execution price is incorporated into
    \(\bar P_{i,j,t}^{\text{entry}}\). If the trade increases the absolute size of an
    existing position without changing direction, the entry price is updated as a
    quantity-weighted average. If the trade closes the position, the entry price is
    cleared. If the trade crosses through zero and opens the opposite direction,
    the entry price is reset to that execution price. This convention gives both
    long and short positions a personal gain/loss anchor while preserving the
    realised execution history.

    The recent peak-wealth state is
    \[
        W_{i,t}^{\max}
        =\max\{W_{i,u}:0\le u\le t,\ W_{i,u}>0\}.
    \]
    It is refreshed whenever the agent's mark-to-market wealth is refreshed and is
    a portfolio quantity, so it responds to gains and losses in every asset. The
    moving-average price reference is asset-specific,
    \[
        \MA_{j,t}^{(h_{\nu})}
        =
        \frac{1}{h_{\nu}}\sum_{u=t-h_{\nu}+1}^{t}P_{j,u},
    \]
    when at least \(h_{\nu}\) market observations are available; otherwise the
    latest market price \(P_{j,t}\) is used for that mode.

    The resolved anchor is therefore
    \[
        \nu_{i,j,t}=
        \begin{cases}
            \nu_i^0,
            & m_{\nu}=\mathrm{configured},\\[0.8ex]
            \bar P_{i,j,t}^{\text{entry}},
            & m_{\nu}=\mathrm{entry}\ \text{and}\ \bar P_{i,j,t}^{\text{entry}}>0,\\[0.8ex]
            W_{i,t}^{\max},
            & m_{\nu}=\mathrm{peakwealth}\ \text{and}\ W_{i,t}^{\max}>0,\\[0.8ex]
            \MA_{j,t}^{(h_{\nu})},
            & m_{\nu}=\mathrm{movingaverage}\ \text{and enough price history is available},\\[0.8ex]
            P_{j,t},
            & m_{\nu}=\mathrm{movingaverage}\ \text{and price history is insufficient},\\[0.8ex]
            W_{i,0},
            & m_{\nu}=\mathrm{initialwealth}\ \text{and}\ W_{i,0}>0,\\[0.8ex]
            \nu_i^0,
            & \text{otherwise}.
        \end{cases}
    \]

    The anchor is then mapped into return space. Price-anchored modes
    (\(\mathrm{configured}\), \(\mathrm{entry}\), \(\mathrm{movingaverage}\)) are
    divided by the asset's current price, while wealth-anchored modes
    (\(\mathrm{peakwealth}\), \(\mathrm{initialwealth}\)) are divided by current
    wealth:
    \[
        \varrho_{i,j,t}
        =
        \clip\!\left(
            \log\frac{\nu_{i,j,t}}{B_{i,j,t}};\,-0.95,0.95
        \right),
    \]
    \[
        B_{i,j,t}=
        \begin{cases}
            W_{i,t}, & m_{\nu}\in\{\mathrm{peakwealth},\mathrm{initialwealth}\},\\
            P_{j,t}, & \text{otherwise},
        \end{cases}
    \]
    with \(\varrho_{i,j,t}=0\) when the base is not a usable positive value. The
    raw anchor \(\nu_{i,j,t}\) is still recorded on the agent for the final
    snapshot, so an experiment can audit which personal anchor was active at the
    last behavioural decision.

    Working in return space is what makes behavioural risk appetite invariant to
    price and wealth levels. An entry price of \(101\) against a price of \(100\)
    and an entry price of \(40.4\) against a price of \(40\) now describe the same
    behavioural state --- a one-percent unrealised loss --- and produce the same
    adjustment, which is the property a multi-asset CPT population requires.

    The default configuration uses \(m_{\nu}=\mathrm{entry}\), corresponding to the
    implemented key value \texttt{average\_purchase\_price}. This choice makes each
    agent evaluate prospective gains and losses relative to its own position
    anchor in that asset, while \texttt{configured} preserves the old
    single-reference baseline for controlled comparison experiments.

    \subsection{Preference Heterogeneity}

    When the global CPT switch is enabled, the implementation first selects
    \(\mathrm{round}(\varphi_{\mathrm{CPT}}N)\) agents using a deterministic stream
    derived from the simulation seed. This allows mixed expected-utility and CPT
    populations while preserving reproducibility when the assignment fraction is
    changed. If prospect-parameter heterogeneity is disabled, all CPT-assigned
    agents share the configured baseline CPT parameters. If heterogeneity is
    enabled, the implementation draws agent-specific preference parameters at
    initialisation:
    \[
        \lambda_i\sim
        \mathrm{TN}\!\left(\lambda,\sigma_\lambda^2;\epsilon_\lambda,10\right),
    \]
    \[
        \alpha_i^G\sim
        \mathrm{TN}\!\left(\alpha^G,\sigma_{\alpha^G}^2;\epsilon_\alpha,1.5\right),
        \qquad
        \alpha_i^L\sim
        \mathrm{TN}\!\left(\alpha^L,\sigma_{\alpha^L}^2;\epsilon_\alpha,1.5\right),
    \]
    \[
        \psi_i^G=\psi_i^L\sim
        \mathcal{U}\!\left(\psi_{\min},\psi_{\max}\right).
    \]
    Here \(\mathrm{TN}(\mu,\sigma^2;a,b)\) denotes a normal draw truncated to
    \([a,b]\). The draw is deterministic conditional on the simulation seed, so a
    given configuration can be reproduced exactly.

    \section{Stylised Facts Measured}\label{sec:stylised_facts}

    The model output is evaluated through stylised-fact diagnostics applied to the
    simulated market-price series. Every enabled asset is screened separately and
    on identical terms, so an \(n\)-asset run reports \(n\) independent verdicts
    rather than one verdict for whichever asset happens to be nominated primary.
    Wealth inequality is a portfolio-level quantity and is therefore repeated
    unchanged on each asset's row. Throughout this section the asset index \(j\)
    is fixed and suppressed where no ambiguity arises, and \(T\) denotes the
    number of analysed return observations.

    The default analysis uses the simulation phase after warm-up and computes
    one-period log returns
    \[
        \ell_{j,t}=\log\frac{P_{j,t}}{P_{j,t-1}},
        \qquad t=1,\ldots,T.
    \]
    Let
    \[
        \bar\ell_j=\frac{1}{T}\sum_{t=1}^{T}\ell_{j,t}
    \]
    denote the sample mean return. The diagnostics focus on three core empirical
    regularities: weak autocorrelation in raw returns, fat-tailed return
    distributions, and volatility clustering. These targets follow the stylised
    facts commonly evaluated in empirical market-microstructure and
    heterogeneous-agent market models.

    \subsection{Absence of Linear Return Autocorrelation}

    For lag \(k\), the sample autocorrelation of raw log returns is
    \[
        \hat\rho_k
        =
        \frac{\sum_{t=k+1}^{T}\left(\ell_{j,t}-\bar\ell_j\right)\left(\ell_{j,t-k}-\bar\ell_j\right)}
        {\sum_{t=1}^{T}{\left(\ell_{j,t}-\bar\ell_j\right)}^2}.
    \]
    The empirical stylised fact is not that every finite-sample autocorrelation is
    exactly zero, but that raw returns show little predictable linear dependence.
    This is assessed using the Ljung-Box statistic at selected horizons
    \(h\in\{10,20,30\}\):
    \[
        Q_{\mathrm{LB}}(h)
        =
        T(T+2)\sum_{k=1}^{h}\frac{\hat\rho_k^2}{T-k}.
    \]
    The null hypothesis is
    \[
        H_0:\rho_1=\rho_2=\cdots=\rho_h=0.
    \]
    A simulated run is interpreted as consistent with this stylised fact when the
    test does not reject \(H_0\) at the chosen significance level for the tested
    lags.

    \subsection{Fat-Tailed Return Distribution}

    Fat tails are measured using the fourth central moment of the simulated
    returns. The sample kurtosis is
    \[
        \widehat K
        =
        \frac{\frac{1}{T}\sum_{t=1}^{T}{\left(\ell_{j,t}-\bar\ell_j\right)}^4}
        {{\left[\frac{1}{T}\sum_{t=1}^{T}{\left(\ell_{j,t}-\bar\ell_j\right)}^2\right]}^2},
    \]
    and excess kurtosis is
    \[
        \widehat K_{\mathrm{ex}}=\widehat K-3.
    \]
    The Gaussian benchmark has \(\widehat K_{\mathrm{ex}}\approx 0\). Positive
    excess kurtosis indicates that large absolute returns occur more frequently
    than under a normal benchmark.

    The Jarque-Bera statistic is also reported:
    \[
        JB
        =
        \frac{T}{6}
        \left(
            \widehat S^2+\frac{\widehat K_{\mathrm{ex}}^2}{4}
        \right),
    \]
    where \(\widehat S\) is sample skewness. The normality null is rejected when the
    Jarque-Bera \(p\)-value is below the significance level. A run is interpreted as
    supporting fat tails when excess kurtosis is positive and normality is rejected.

    \subsection{Volatility Clustering}

    Volatility clustering means that large price movements tend to be followed by
    large movements, and small movements by small movements, even when the signs of
    returns remain difficult to predict. The model therefore evaluates dependence in
    absolute and squared returns:
    \[
        \hat\rho_k^{|\ell|}
        =
        \mathrm{Corr}(|\ell_{j,t}|,|\ell_{j,t-k}|),
        \qquad
        \hat\rho_k^{\ell^2}
        =
        \mathrm{Corr}(\ell_{j,t}^2,\ell_{j,t-k}^2).
    \]
    Positive persistence in these autocorrelation functions is visual evidence of
    clustered volatility.

    The formal test uses ARCH effects. For a chosen lag order \(q\), define squared
    demeaned returns
    \[
        u_t^2={\left(\ell_{j,t}-\bar\ell_j\right)}^2.
    \]
    The auxiliary regression is
    \[
        u_t^2
        =
        a_0+\sum_{j=1}^{q}a_j u_{t-j}^2+\varepsilon_t.
    \]
    The ARCH LM statistic is
    \[
        LM_q=TR_q^2,
    \]
    where \(R_q^2\) is the coefficient of determination from the auxiliary
    regression. The null hypothesis is no ARCH effect:
    \[
        H_0:a_1=a_2=\cdots=a_q=0.
    \]
    Rejection of this null supports volatility clustering.

    When the required package is available, the diagnostics also estimate a
    GARCH \((1,1)\) volatility equation
    \[
        \sigma_t^2
        =
        \omega+\alpha_G u_{t-1}^2+\beta_G\sigma_{t-1}^2.
    \]
    The persistence measure is
    \[
        \alpha_G+\beta_G.
    \]
    Values close to one indicate highly persistent conditional volatility, which is
    typical of crisis-sensitive financial return series.

    \section{Reproducible Experiment Identifier}\label{sec:experiment_identifier}

    When experiment output archiving is enabled, the run folder records the main
    experimental controls in numeric form:
    \begin{center}
        \texttt{YYYYMMDDHHMMSS\_prospect\_theory\_prospect\_heterogeneity\_}\\
        \texttt{social\_enabled\_short\_sale\_num\_fundamentalists\_num\_chartists\_}\\
        \texttt{num\_noise\_traders\_num\_buy\_and\_hold\_agents\_time\_steps\_seed}.
    \end{center}
    Boolean switches are encoded as \(1\) when enabled and \(0\) when disabled. For
    example,
    \begin{center}
        \texttt{20260427143012\_1\_1\_0\_1\_100\_1500\_25\_1\_1000\_66}
    \end{center}
    denotes a run created at local timestamp \(2026\)-\(04\)-\(27\) \(14{:}30{:}12\),
    with CPT enabled, heterogeneous CPT parameters enabled, social influence
    disabled, short selling enabled, \(100\) fundamentalists, \(1500\) chartists,
    \(25\) noise traders, \(1\) passive buy-and-hold agent, \(1000\) simulation
    steps, and seed \(66\).

    The identifier schema is deliberately unchanged by the multi-asset extension.
    Asset composition is recorded inside the run folder --- in the resolved
    configuration and metadata, which also list assets that were declared but
    disabled --- rather than in the folder name, so run identifiers stay
    comparable across single-asset and multi-asset experiments.

    \section{Multi-Asset Output Convention}\label{sec:multi_asset_outputs}

    Two views of the same run are exported, and they are not alternatives: the
    long-format view is the analytical one, and the single-asset view exists for
    continuity.

    The long-format per-asset time series contains one row per
    \((\text{phase},\text{step},\text{asset})\) with that asset's market and
    fundamental price, simple and log return, submitted limit and market order
    counts, traded volume, executed trades, last trade price, best bid and ask,
    spread, midquote, bid and ask depth, resting order counts, the five
    cancellation counters of Section~\ref{sec:continuous_lob} (stale, expired,
    infeasible-resting, infeasible-incoming, self-trade prevention) together with
    collar rejections, transaction costs, the allocative-efficiency triple
    \((AS_{j,t},MS_{j,t},AE_{j,t})\), the mean agent portfolio weight
    \(\bar w_{j,t}\), and the count of agents entering that mean. Initial and
    warm-up rows carry prices only; the trading-phase fields are left empty for
    them.

    The single-asset step time series retains its original schema and reports the
    primary asset \(j^{\ast}\) only. It exists so that calibration and legacy
    tooling continue to run unchanged. Any conclusion about the market as a whole
    should be drawn from the per-asset view and the per-asset stylised-fact
    verdicts, since the primary asset is a reporting projection and not a
    privileged market.

\end{document}
